% Generated mechanically from frozen input inputs/p10/ja/spaces-duality.tex.
% Do not edit this generated file; edit the bound locale source instead.

% OK, start here.
%

\setcounter{chapter}{85}
\chapter{空間の双対性}



\phantomsection
\label{spaces-duality-section-phantom}


\section{序論}
\label{spaces-duality-section-introduction}

\noindent
本章は、スキームに対する対応する章の類似物である。『スキームの双対性』節
\ref{duality-section-introduction} を参照されたい。
ここでの展開は、論文
\cite{Neeman-Grothendieck}, \cite{LN},
\cite{Lipman-notes}, \cite{Neeman-improvement}
における展開と同様である。




\section{代数空間上の双対化複体}
\label{spaces-duality-section-dualizing-spaces}

\noindent
$U$ を局所 Noether スキームとする。$U$ の小エタール・サイト上の
$U$ の構造層を $\mathcal{O}_\etale$ とする。
『スキームの双対性』節
\ref{duality-section-dualizing-schemes} の意味でのある双対化複体
$\omega_U^\bullet$ に対して $K = \epsilon^*(\omega_U^\bullet)$
と書けるとき、対象 $K \in D_\QCoh(\mathcal{O}_\etale)$ を
$U$ 上の双対化複体という。ここで
$\epsilon^* : D_\QCoh(\mathcal{O}_U) \to D_\QCoh(\mathcal{O}_\etale)$
は『空間の導来圏』補題
\ref{spaces-perfect-lemma-derived-quasi-coherent-small-etale-site}
の同値である。『スキームの双対性』節
\ref{duality-section-dualizing-schemes} で調べた
$\omega_U^\bullet$ の性質の大部分は、『空間の導来圏』節
\ref{spaces-perfect-section-derived-quasi-coherent-etale} および
\ref{spaces-perfect-section-spell-out} の議論を通じて $K$ に受け継がれる。

\medskip\noindent
局所 Noether 代数空間上の双対化複体を、エタール局所的に
対応するスキーム上の双対化複体から来る複体として定義する。

\begin{lemma}
\label{spaces-duality-lemma-equivalent-definitions}
$S$ をスキームとし、$X$ を $S$ 上の局所 Noether 代数空間とする。
$K$ を $D_\QCoh(\mathcal{O}_X)$ の対象とする。次の条件は同値である。
\begin{enumerate}
\item $U$ がスキームである任意のエタール射 $U \to X$ に対して、
制限 $K|_U$ は $U$ 上の双対化複体である（意味は上記のとおり）。
\item $U$ がスキームである全射エタール射 $U \to X$ で、
$K|_U$ が $U$ 上の双対化複体となるものが存在する。
\end{enumerate}
\end{lemma}

\begin{proof}
$U$ がスキームである全射エタール射 $U \to X$ を仮定する。
$V$ がスキームであるエタール射 $V \to X$ をとる。このとき
$$
U \leftarrow U \times_X V \rightarrow V
$$
はスキームのエタール射であり、$V$ への射は全射である。
したがって『スキームの双対性』補題
\ref{duality-lemma-descent-ascent} により、$K|_U$ が $U$ 上の
双対化複体ならば、$K|_V$ は $V$ 上の双対化複体である。
\end{proof}

\begin{definition}
\label{spaces-duality-definition-dualizing-scheme}
$S$ をスキームとする。
$X$ を $S$ 上の局所 Noether 代数空間とする。
$D_\QCoh(\mathcal{O}_X)$ の対象 $K$ が補題
\ref{spaces-duality-lemma-equivalent-definitions} の同値な条件を満たすとき、
$K$ を {\it 双対化複体}という。
\end{definition}

\begin{lemma}
\label{spaces-duality-lemma-affine-duality}
$A$ を Noether 環とし、$X = \Spec(A)$ とする。
$X$ の小エタール・サイト上の $X$ の構造層を
$\mathcal{O}_\etale$ とする。$K, L$ を $D(A)$ の対象とする。
$K \in D_{\textit{Coh}}(A)$ であり、かつ $L$ の入射次元が有限ならば、
$$
\epsilon^*\widetilde{R\Hom_A(K, L)} =
R\SheafHom_{\mathcal{O}_\etale}(\epsilon^*\widetilde{K},
\epsilon^*\widetilde{L})
$$
が $D(\mathcal{O}_\etale)$ において成り立つ。ここで
$\epsilon : (X_\etale, \mathcal{O}_\etale) \to (X, \mathcal{O}_X)$
は『空間の導来圏』節
\ref{spaces-perfect-section-derived-quasi-coherent-etale}
におけるものとする。
\end{lemma}

\begin{proof}
『スキームの双対性』補題 \ref{duality-lemma-affine-duality} により、
標準同型
$$
\widetilde{R\Hom_A(K, L)} =
R\SheafHom_{\mathcal{O}_X}(\widetilde{K}, \widetilde{L})
$$
が $D(\mathcal{O}_X)$ において存在する。また、標準射
$$
\epsilon^*R\Hom_{\mathcal{O}_X}(\widetilde{K}, \widetilde{L})
\longrightarrow
R\SheafHom_{\mathcal{O}_\etale}(\epsilon^*\widetilde{K},
\epsilon^*\widetilde{L})
$$
が $D(\mathcal{O}_\etale)$ において存在する。『サイト上のコホモロジー』評注
\ref{sites-cohomology-remark-prepare-fancy-base-change} を参照されたい。
この射の左辺と右辺のコホモロジー層が同型であることを示すが、
その同型がこの射によって与えられることの検証は省略する。

\medskip\noindent
$L$ は入射 $A$-加群からなる有限複体 $I^\bullet$ で与えられると
仮定してよい。$I^\bullet$ の長さに関する帰納法と、各構成の
識別三角形との両立性により、$I$ が入射 $A$-加群で
$L = I[0]$ である場合に帰着する。
% P10-E0246
$R\SheafHom_{\mathcal{O}_\etale}(\epsilon^*\widetilde{K},
\epsilon^*\widetilde{L}))$ のコホモロジー層は、$X$ 上エタールな
$U$ に対し、$U_\etale$ への $\epsilon^*\widetilde{K}$ と
$\epsilon^*\widetilde{L}$ の制限の間の第 $i$ Ext 群を対応させる
前層の層化であることを思い出そう。『サイト上のコホモロジー』補題
\ref{sites-cohomology-lemma-section-RHom-over-U} を参照されたい。
$U = \Spec(B)$ がアフィンなら、この Ext 群は、『空間の導来圏』補題
\ref{spaces-perfect-lemma-derived-quasi-coherent-small-etale-site} と
『スキームの導来圏』補題
\ref{perfect-lemma-affine-compare-bounded} の同値により、
$\text{Ext}^i_B(K \otimes_A B, L \otimes_A B)$ に等しい
（ここでは『空間の導来圏』評注
\ref{spaces-perfect-remark-match-total-direct-images}
に詳述された両立性も用いる）。$A \to B$ はエタールなので、
『双対化複体』補題 \ref{dualizing-lemma-injective-goes-up} により、
$I \otimes_A B$ は入射 $B$-加群である。したがって
\begin{align*}
\Ext^n_B(K \otimes_A B, I \otimes_A B)
& =
\Hom_B(H^{-n}(K \otimes_A B), I \otimes_A B) \\
% P10-E0247
& =
\Hom_{A_f}(H^{-n}(K) \otimes_A B, I \otimes_A B) \\
& =
\Hom_A(H^{-n}(K), I) \otimes_A B \\
& =
\text{Ext}^n_A(K, I) \otimes_A B
\end{align*}
を得る。最後から二つ目の等号は、$H^{-n}(K)$ が有限 $A$-加群であることによる。
% P10-E0248
『代数詳論』補題
\ref{more-algebra-lemma-pseudo-coherence-and-base-change-ext}
を参照されたい。ゆえに、補題の等式の左辺と右辺のコホモロジー層は同じである。
\end{proof}

\begin{lemma}
\label{spaces-duality-lemma-dualizing-spaces}
$S$ をスキームとし、$X$ を $S$ 上の局所 Noether 代数空間とする。
$K$ を $X$ 上の双対化複体とする。このとき
$K$ は $D_{\textit{Coh}}(\mathcal{O}_X)$ の対象であり、
$D = R\SheafHom_{\mathcal{O}_X}(-, K)$ は反同値
$$
D :
D_{\textit{Coh}}(\mathcal{O}_X)
\longrightarrow
D_{\textit{Coh}}(\mathcal{O}_X)
$$
を誘導する。この反同値には標準同型
$\text{id} \to D \circ D$ が備わる。$X$ が準コンパクトなら、
$D$ は $D^+_{\textit{Coh}}(\mathcal{O}_X)$ と
$D^-_{\textit{Coh}}(\mathcal{O}_X)$ を交換し、同値
$D^b_{\textit{Coh}}(\mathcal{O}_X) \to D^b_{\textit{Coh}}(\mathcal{O}_X)$
を誘導する。
\end{lemma}

\begin{proof}
$U$ がアフィンであるエタール射 $U \to X$ をとる。
$U = \Spec(A)$ と書き、補題 \ref{spaces-duality-lemma-equivalent-definitions} および
『スキームの双対性』補題
\ref{duality-lemma-equivalent-definitions} の意味で $K|_U$ に対応する
$A$ の双対化複体を $\omega_A^\bullet$ とする。
補題 \ref{spaces-duality-lemma-affine-duality} により、図式
$$
\xymatrix{
D_{\textit{Coh}}(A) \ar[r] \ar[d]_{R\Hom_A(-, \omega_A^\bullet)} &
D_{\textit{Coh}}(\mathcal{O}_\etale)
\ar[d]^{R\SheafHom_{\mathcal{O}_\etale}(-, K|_U)} \\
D_{\textit{Coh}}(A) \ar[r] &
D(\mathcal{O}_\etale)
}
$$
は可換である。ここで $\mathcal{O}_\etale$ は $U$ の
小エタール・サイト上の構造層である。$R\SheafHom$ の形成は
制限と可換なので、$D$ は $D_{\textit{Coh}}(\mathcal{O}_X)$ を
$D_{\textit{Coh}}(\mathcal{O}_X)$ に送る。さらに、標準射
$$
L \longrightarrow
R\SheafHom_{\mathcal{O}_X}(R\SheafHom_{\mathcal{O}_X}(L, K), K)
$$
（『サイト上のコホモロジー』補題
\ref{sites-cohomology-lemma-internal-hom-evaluate}）は、
$D_{\textit{Coh}}(\mathcal{O}_X)$ の任意の $L$ に対して同型である。
実際、上記のすべての $U$ 上でこれは『双対化複体』補題
\ref{dualizing-lemma-dualizing} により成り立つ。
準コンパクトな場合の関手 $D$ の有界性に関する主張も、
『双対化複体』補題 \ref{dualizing-lemma-dualizing}
の対応する主張から従う。
\end{proof}

\noindent
$(\mathcal{C}, \mathcal{O})$ を環付きサイトとする。
$D(\mathcal{O})$ の対象 $L$ が {\it 可逆}であるとは、導来テンソル積が
% P10-E0249
$D(\mathcal{O}_X)$ に与える対称モノイド構造に関して
その対象が可逆対象であることをいう。『サイト上のコホモロジー』補題
\ref{sites-cohomology-lemma-invertible-derived} で見たように、これは
$L$ が完全であることを意味する。さらに $(\mathcal{C}, \mathcal{O})$
% P10-E0250
が局所環付きサイトなら、$\mathcal{C}$ の任意の対象 $U$ に対し、
ある整数 $n_i$ について $L|_{U_i} \cong \mathcal{O}_{U_i}[-n_i]$
となるような $\mathcal{C}$ における $U$ の被覆
$\{U_i \to U\}$ が存在する。

\medskip\noindent
$S$ をスキームとし、$X$ を $S$ 上の代数空間とする。
$D(\mathcal{O}_X)$ の $L$ が可逆なら、非交和分解
$X = \coprod_{n \in \mathbf{Z}} X_n$ であって、
$L|_{X_n}$ が次数 $n$ に置かれた可逆加群となるものが存在する。
特に $L = \bigoplus H^n(L)[-n]$ であり、これは $L$ を表す
$\mathcal{O}_X$-加群の明確に定まった複体（微分は零）を与える。

\begin{lemma}
\label{spaces-duality-lemma-dualizing-unique-spaces}
$S$ をスキームとする。
$X$ を $S$ 上の局所 Noether 代数空間とする。
$K$ と $K'$ が $X$ 上の双対化複体なら、$K'$ は、
$D(\mathcal{O}_X)$ のある可逆対象 $L$ に対する
$K \otimes_{\mathcal{O}_X}^\mathbf{L} L$ と同型である。
\end{lemma}

\begin{proof}
$$
L = R\SheafHom_{\mathcal{O}_X}(K, K')
$$
とおく。これは $D(\mathcal{O}_X)$ の可逆対象である。実際、
このことはアフィン局所的に成り立つ。補題
\ref{spaces-duality-lemma-affine-duality} と『双対化複体』補題
\ref{dualizing-lemma-dualizing-unique} およびその証明を用いよ。
同じ理由により、評価写像
$L \otimes_{\mathcal{O}_X}^\mathbf{L} K \to K'$ は同型である。
\end{proof}

\begin{lemma}
\label{spaces-duality-lemma-dimension-function-scheme}
$S$ をスキームとし、$X$ を $S$ 上の局所 Noether かつ
準分離な代数空間とする。$\omega_X^\bullet$ を $X$ 上の双対化複体とする。
% P10-E0251
このとき $X$ について、次で定義される関数 $|X| \to \mathbf{Z}$
$$
\begin{aligned}
x &\longmapsto \delta(x)\text{ であって } \\
&\omega_{X, \overline{x}}^\bullet[-\delta(x)] \\
&\text{ は次の環上の正規化双対化複体である：} \\
&\mathcal{O}_{X, \overline{x}}
\end{aligned}
$$
は $|X|$ 上の次元関数である。
\end{lemma}

\begin{proof}
$U$ をスキームとし、$U \to X$ を全射エタール射とする。
$\omega_X^\bullet|_U$ に付随する $U$ 上の双対化複体を
$\omega_U^\bullet$ とする。$u \in U$ が $x \in |X|$ に写るなら、
$\mathcal{O}_{X, \overline{x}}$ は $\mathcal{O}_{U, u}$ の強 Hensel 化である。
『双対化複体』補題 \ref{dualizing-lemma-flat-unramified} により、
$\omega^\bullet$ が $\mathcal{O}_{U, u}$ の正規化双対化複体なら、
$\omega^\bullet \otimes_{\mathcal{O}_{U, u}}
\mathcal{O}_{X, \overline{x}}$ は
$\mathcal{O}_{X, \overline{x}}$ の正規化双対化複体である。
したがって、スキーム $U$ と複体 $\omega_U^\bullet$ に対する
『スキームの双対性』補題
\ref{duality-lemma-dimension-function-scheme} の次元関数
$U \to \mathbf{Z}$ は、$U \to |X|$ と $\delta$ の合成に等しい。
$|X|$ における特殊化は $U$ における特殊化へ持ち上がり、
$U$ における非自明な特殊化は $X$ における非自明な特殊化へ写る
（『良質な空間』補題
\ref{decent-spaces-lemma-decent-specialization} および
\ref{decent-spaces-lemma-decent-no-specializations-map-to-same-point}）
ことを用いると、容易な位相的議論により $\delta$ は
$|X|$ 上の次元関数であることが分かる。
\end{proof}




\section{順像の右随伴}
\label{spaces-duality-section-twisted-inverse-image}

\noindent
これは『スキームの双対性』節
\ref{duality-section-twisted-inverse-image} の類似物である。

\begin{lemma}
\label{spaces-duality-lemma-twisted-inverse-image}
\begin{reference}
これは \cite[Example 4.2]{Neeman-Grothendieck} とほぼ同じである。
\end{reference}
$S$ をスキームとする。
$f : X \to Y$ を、$S$ 上の準分離かつ準コンパクトな代数空間の間の
射とする。関手 $Rf_* : D_\QCoh(X) \to D_\QCoh(Y)$ は右随伴をもつ。
\end{lemma}

\begin{proof}
『導来圏』命題 \ref{derived-proposition-brown} の仮定を検証することにより、
右随伴が存在することを証明する。まず、圏
$D_\QCoh(\mathcal{O}_X)$ は直和をもつ。『空間の導来圏』補題
\ref{spaces-perfect-lemma-quasi-coherence-direct-sums} を参照されたい。
また『空間の導来圏』定理
\ref{spaces-perfect-theorem-bondal-van-den-Bergh} により、圏
$D_\QCoh(\mathcal{O}_X)$ はコンパクト生成される。
$X$ と $Y$ は準コンパクトかつ準分離なので、$f$ もそうである。
『空間の射』補題
\ref{spaces-morphisms-lemma-compose-after-separated} および
\ref{spaces-morphisms-lemma-quasi-compact-permanence} を参照されたい。
したがって関手 $Rf_*$ は直和と可換である。『空間の導来圏』補題
\ref{spaces-perfect-lemma-quasi-coherence-pushforward-direct-sums}
を参照されたい。これで証明は完了する。
\end{proof}

\begin{lemma}
\label{spaces-duality-lemma-twisted-inverse-image-bounded-below}
記法と仮定は補題 \ref{spaces-duality-lemma-twisted-inverse-image} のとおりとする。
$a : D_\QCoh(\mathcal{O}_Y) \to D_\QCoh(\mathcal{O}_X)$ を
$Rf_*$ の右随伴とする。このとき $a$ は
$D^+_\QCoh(\mathcal{O}_Y)$ を $D^+_\QCoh(\mathcal{O}_X)$ に写す。
実際、ある整数 $N$ が存在し、$i \leq c$ に対して
$H^i(K) = 0$ ならば、$i \leq c - N$ に対して
$H^i(a(K)) = 0$ となる。
\end{lemma}

\begin{proof}
『空間の導来圏』補題
\ref{spaces-perfect-lemma-quasi-coherence-direct-image} により、
関手 $Rf_*$ のコホモロジー次元は有限である。言い換えると、
ある整数 $N$ が存在し、$i \geq c$ に対して $H^i(L) = 0$ ならば、
$i \geq N + c$ に対して $H^i(Rf_*L) = 0$ となる。
$K \in D^+_\QCoh(\mathcal{O}_Y)$ が $i \leq c$ に対して
$H^i(K) = 0$ を満たすとする。このとき、上で述べたことにより
$$
\Hom_{D(\mathcal{O}_X)}(\tau_{\leq c - N}a(K), a(K)) =
\Hom_{D(\mathcal{O}_Y)}(Rf_*\tau_{\leq c - N}a(K), K) = 0
$$
である。明らかに、これは $i \leq c - N$ に対して
$H^i(a(K)) = 0$ であることを含意する。
\end{proof}

\noindent
$S$ をスキームとする。
$f : X \to Y$ を、$S$ 上の準分離かつ準コンパクトな代数空間の射とする。
$Rf_* : D_\QCoh(\mathcal{O}_X) \to D_\QCoh(\mathcal{O}_Y)$
の右随伴を $a$ と書く。任意の
$K \in D_\QCoh(\mathcal{O}_Y)$ と
$L \in D_\QCoh(\mathcal{O}_X)$ に対して、標準射
\begin{equation}
\label{spaces-duality-equation-sheafy-trace}
Rf_*R\SheafHom_{\mathcal{O}_X}(L, a(K))
\longrightarrow
R\SheafHom_{\mathcal{O}_Y}(Rf_*L, K)
\end{equation}
を得る。すなわち、この射は合成
$$
\begin{aligned}
Rf_*R\SheafHom_{\mathcal{O}_X}(L, a(K))
&\to R\SheafHom_{\mathcal{O}_Y}(Rf_*L, Rf_*a(K)) \\
&\to R\SheafHom_{\mathcal{O}_Y}(Rf_*L, K)
\end{aligned}
$$
として構成される。ここで最初の矢印は『サイト上のコホモロジー』評注
\ref{sites-cohomology-remark-projection-formula-for-internal-hom}
のものであり、二番目の矢印は随伴の余単位 $Rf_*a(K) \to K$ である。

\begin{lemma}
\label{spaces-duality-lemma-iso-on-RSheafHom}
$S$ をスキームとする。
$f : X \to Y$ を、$S$ 上の準コンパクトかつ準分離な代数空間の射とする。
$a$ を $Rf_* : D_\QCoh(\mathcal{O}_X) \to
D_\QCoh(\mathcal{O}_Y)$ の右随伴とする。
$L \in D_\QCoh(\mathcal{O}_X)$ および
$K \in D_\QCoh(\mathcal{O}_Y)$ とする。このとき、写像
(\ref{spaces-duality-equation-sheafy-trace})
$$
Rf_*R\SheafHom_{\mathcal{O}_X}(L, a(K))
\longrightarrow
R\SheafHom_{\mathcal{O}_Y}(Rf_*L, K)
$$
は、『空間の導来圏』節
\ref{spaces-perfect-section-better-coherator} で論じた関手
$DQ_Y : D(\mathcal{O}_Y) \to D_\QCoh(\mathcal{O}_Y)$
を適用すると同型になる。
\end{lemma}

\begin{proof}
『空間の導来圏』補題
\ref{spaces-perfect-lemma-better-coherator} により $DQ_Y$ が存在するので、
この主張には意味がある。$DQ_Y$ は包含関手
$D_\QCoh(\mathcal{O}_Y) \to D(\mathcal{O}_Y)$ の右随伴なので、
補題を証明するには、任意の $M \in D_\QCoh(\mathcal{O}_Y)$ に対して、
% P10-E0252
写像 (\ref{spaces-duality-equation-sheafy-trace}) が全単射
$$
\Hom_Y(M, Rf_*R\SheafHom_{\mathcal{O}_X}(L, a(K)))
\longrightarrow
\Hom_Y(M, R\SheafHom_{\mathcal{O}_Y}(Rf_*L, K))
$$
を誘導することを示せばよい。このために、次の等式列を用いる：
\begin{align*}
\Hom_Y(M, Rf_*R\SheafHom_{\mathcal{O}_X}(L, a(K)))
& =
\Hom_X(Lf^*M, R\SheafHom_{\mathcal{O}_X}(L, a(K))) \\
& =
\Hom_X(Lf^*M \otimes_{\mathcal{O}_X}^\mathbf{L} L, a(K)) \\
& =
\Hom_Y(Rf_*(Lf^*M \otimes_{\mathcal{O}_X}^\mathbf{L} L), K) \\
& =
\Hom_Y(M \otimes_{\mathcal{O}_Y}^\mathbf{L} Rf_*L, K) \\
& =
\Hom_Y(M, R\SheafHom_{\mathcal{O}_Y}(Rf_*L, K))
\end{align*}
最初の等号は『サイト上のコホモロジー』補題
\ref{sites-cohomology-lemma-adjoint} による。
% P10-E0253
二番目の等号は『サイト上のコホモロジー』補題
\ref{sites-cohomology-lemma-internal-hom} による。
三番目の等号は $a$ の構成による。
四番目の等号は『空間の導来圏』補題
\ref{spaces-perfect-lemma-cohomology-base-change} による
（これが重要な段階である）。五番目の等号は『サイト上のコホモロジー』補題
\ref{sites-cohomology-lemma-internal-hom} による。
\end{proof}

\begin{example}
\label{spaces-duality-example-iso-on-RSheafHom}
補題 \ref{spaces-duality-lemma-iso-on-RSheafHom} の主張は、「コヒーレータ」$DQ_Y$
を適用しなければ成り立たない。『スキームの双対性』例
\ref{duality-example-iso-on-RSheafHom} を参照されたい。
\end{example}

\begin{remark}
\label{spaces-duality-remark-iso-on-RSheafHom}
補題 \ref{spaces-duality-lemma-iso-on-RSheafHom} の状況では、『空間の導来圏』補題
\ref{spaces-perfect-lemma-pushforward-better-coherator} により
$$
DQ_Y(Rf_*R\SheafHom_{\mathcal{O}_X}(L, a(K))) =
Rf_* DQ_X(R\SheafHom_{\mathcal{O}_X}(L, a(K)))
$$
である。したがって
$R\SheafHom_{\mathcal{O}_X}(L, a(K))\allowbreak \in D_\QCoh(\mathcal{O}_X)$
なら、矢印の左辺にある $DQ_Y$ を「消去」できる。一方、
$R\SheafHom_{\mathcal{O}_Y}(Rf_*L, K)\allowbreak \in D_\QCoh(\mathcal{O}_Y)$
であることが分かっていれば、矢印の右辺から $DQ_Y$ を「消去」できる。
両方が成り立つなら、(\ref{spaces-duality-equation-sheafy-trace}) は同型である。
これと『空間の導来圏』補題
\ref{spaces-perfect-lemma-quasi-coherence-internal-hom} を組み合わせると、
次のいずれかの場合に
$Rf_*R\SheafHom_{\mathcal{O}_X}(L, a(K))\allowbreak \to
R\SheafHom_{\mathcal{O}_Y}(Rf_*L, K)$ は同型であることが分かる。
\begin{enumerate}
\item $L$ と $Rf_*L$ が完全である。
\item $K$ が下に有界であり、$L$ と $Rf_*L$ が擬連接である。
\end{enumerate}
(2) では、補題 \ref{spaces-duality-lemma-twisted-inverse-image-bounded-below} により、
$K$ が下に有界なら $a(K)$ も下に有界であることを用いる。
\end{remark}

\begin{example}
\label{spaces-duality-example-iso-on-RSheafHom-noetherian}
$S$ をスキームとする。
$f : X \to Y$ を $S$ 上の Noether 代数空間の固有射とし、
$L \in D^-_{\textit{Coh}}(X)$ および
$K \in D^+_{\QCoh}(\mathcal{O}_Y)$ とする。このとき写像
$Rf_*R\SheafHom_{\mathcal{O}_X}(L, a(K)) \to
R\SheafHom_{\mathcal{O}_Y}(Rf_*L, K)$ は同型である。
実際、『空間の導来圏』補題
\ref{spaces-perfect-lemma-identify-pseudo-coherent-noetherian} および
\ref{spaces-perfect-lemma-direct-image-coherent} により、複体 $L$ と
$Rf_*L$ は擬連接であり、評注 \ref{spaces-duality-remark-iso-on-RSheafHom}
の議論を適用できる。
\end{example}

\begin{lemma}
\label{spaces-duality-lemma-iso-global-hom}
$S$ をスキームとする。
$f : X \to Y$ を、$S$ 上の準分離かつ準コンパクトな代数空間の射とする。
すべての $L \in D_\QCoh(\mathcal{O}_X)$ と
$K \in D_\QCoh(\mathcal{O}_Y)$ に対して、
(\ref{spaces-duality-equation-sheafy-trace}) は大域導来 Hom の同型
$R\Hom_X(L, a(K))\allowbreak \to R\Hom_Y(Rf_*L, K)$ を誘導する。
\end{lemma}

\begin{proof}
構成により（『サイト上のコホモロジー』節
\ref{sites-cohomology-section-global-RHom}）、複体
$$
\begin{aligned}
R\Hom_X(L, a(K))
&= R\Gamma(X, R\SheafHom_{\mathcal{O}_X}(L, a(K))) \\
&= R\Gamma(Y, Rf_*R\SheafHom_{\mathcal{O}_X}(L, a(K)))
\end{aligned}
$$
および
$$
% P10-E0254
R\Hom_Y(Rf_*L, K) = R\Gamma(Y, R\SheafHom_{\mathcal{O}_X}(Rf_*L, a(K)))
$$
を得る。したがって補題は補題 \ref{spaces-duality-lemma-iso-on-RSheafHom} の帰結である。
すなわち、$D(\mathcal{O}_Y)$ の写像 $E \to E'$ が同型
$DQ_Y(E) \to DQ_Y(E')$ を誘導するなら、大域切断上で準同型
$R\Gamma(Y, E) \to R\Gamma(Y, E')$ を誘導する。実際、
$\mathcal{O}_Y[-i]$ は $D_\QCoh(\mathcal{O}_Y)$ に属し、$DQ_Y$ は
包含関手 $D_\QCoh(\mathcal{O}_Y) \to D(\mathcal{O}_Y)$ の
右随伴である。したがって
\par\noindent
$\begin{aligned}
&H^i(Y, E) = \Ext^i_Y(\mathcal{O}_Y, E) \\
&= \Hom(\mathcal{O}_Y[-i], E) \\
&= \Hom(\mathcal{O}_Y[-i], DQ_Y(E))
\end{aligned}$
である。
\end{proof}






\section{順像の右随伴と基底変換 I}
\label{spaces-duality-section-base-change-map}

\noindent
順像の右随伴の間の基底変換写像を定義しよう。
$S$ をスキームとする。Cartesian 図式
\begin{equation}
\label{spaces-duality-equation-base-change}
\vcenter{
\xymatrix{
X' \ar[r]_{g'} \ar[d]_{f'} & X \ar[d]^f \\
Y' \ar[r]^g & Y
}
}
\end{equation}
を考える。ここで $Y'$ と $X$ は $Y$ 上 {\bf Tor 独立}である。
$Rf_*$ と $Rf'_*$ の右随伴（補題
\ref{spaces-duality-lemma-twisted-inverse-image}）をそれぞれ
$$
a  : D_\QCoh(\mathcal{O}_Y) \to D_\QCoh(\mathcal{O}_X)
\quad\text{および}\quad
a' : D_\QCoh(\mathcal{O}_{Y'}) \to D_\QCoh(\mathcal{O}_{X'})
$$
と書く。『サイト上のコホモロジー』評注
\ref{sites-cohomology-remark-base-change} の基底変換写像は、
準連接コホモロジーをもつ層の導来圏上の関手の変換
$$
Lg^* \circ Rf_* \longrightarrow Rf'_* \circ L(g')^*
$$
を与える。したがって、反対向きの右随伴の間の変換
$$
a \circ Rg_* \longleftarrow Rg'_* \circ a'
$$
を得る。

\begin{lemma}
\label{spaces-duality-lemma-flat-precompose-pus}
図式 (\ref{spaces-duality-equation-base-change}) において、写像
$a \circ Rg_* \leftarrow Rg'_* \circ a'$ は同型である。
\end{lemma}

\begin{proof}
『空間の導来圏』補題
\ref{spaces-perfect-lemma-compare-base-change} により、
$D_\QCoh(\mathcal{O}_X)$ の任意の $K$ に対して、基底変換写像
$Lg^* \circ Rf_* K \to Rf'_* \circ L(g')^*K$ は同型である
（ここで Tor 独立性の仮定を用いる）。したがって、対応する随伴関手間の
変換も同型である。
\end{proof}

\noindent
そこで、次の合成で与えられる関手
$D_\QCoh(\mathcal{O}_Y) \to D_\QCoh(\mathcal{O}_{X'})$
の射を考えることができる：
\begin{equation}
\label{spaces-duality-equation-base-change-map}
L(g')^* \circ a \to
L(g')^* \circ a \circ Rg_* \circ Lg^* \leftarrow
L(g')^* \circ Rg'_* \circ a' \circ Lg^* \to a' \circ Lg^*.
\end{equation}
最初の矢印は随伴写像 $\text{id} \to Rg_* Lg^*$ から来て、
最後の矢印は随伴写像 $L(g')^*Rg'_* \to \text{id}$ から来る。
中央の矢印を逆にするには Tor 独立性の仮定が必要である。
補題 \ref{spaces-duality-lemma-flat-precompose-pus} を参照されたい。
別の見方をすると、$L(g')^*$ と $R(g')_*$ の随伴性により、
(\ref{spaces-duality-equation-base-change-map}) は自然変換
$$
a \to a \circ Rg_* \circ Lg^* \leftarrow Rg'_* \circ a' \circ Lg^*
$$
と考えられる。ここでも二番目の矢印は可逆である。% P10-E0255
$M \in D_\QCoh(\mathcal{O}_X)$ および
$K \in D_\QCoh(\mathcal{O}_Y)$ なら、Yoneda 関手上でこの写像は
\begin{align*}
\Hom_X(M, a(K))
& =
\Hom_Y(Rf_*M, K) \\
& \to
\Hom_Y(Rf_*M, Rg_* Lg^*K) \\
& =
\Hom_{Y'}(Lg^*Rf_*M, Lg^*K) \\
& \leftarrow
\Hom_{Y'}(Rf'_* L(g')^*M, Lg^*K) \\
& =
\Hom_{X'}(L(g')^*M, a'(Lg^*K)) \\
& =
\Hom_X(M, Rg'_*a'(Lg^*K))
\end{align*}
で与えられる（左向きの矢印は『空間の導来圏』補題
\ref{spaces-perfect-lemma-compare-base-change}
の基底変換定理により可逆である）。% P10-E0256
この表示の方が多少明示的である。

\medskip\noindent
本節ではまず、この基底変換写像が、通常の基底変換写像についての
『サイト上のコホモロジー』評注
\ref{sites-cohomology-remark-compose-base-change} および
\ref{sites-cohomology-remark-compose-base-change-horizontal}
と同様に、正方形を積み重ねることに関するいくつかの自然な両立性を
満たすことを証明する。初読では本節の残りを飛ばすことを勧める。

\begin{lemma}
\label{spaces-duality-lemma-compose-base-change-maps}
$S$ をスキームとする。$S$ 上の準コンパクトかつ準分離な
代数空間の可換図式
$$
\xymatrix{
X' \ar[r]_k \ar[d]_{f'} & X \ar[d]^f \\
Y' \ar[r]^l \ar[d]_{g'} & Y \ar[d]^g \\
Z' \ar[r]^m & Z
}
$$
を考える。両方の正方形は Cartesian であり、$f$ と $l$、
および $g$ と $m$ は Tor 独立であるとする。このとき、
二つの正方形に対する写像 (\ref{spaces-duality-equation-base-change-map}) の合成は、
外側の長方形に対する基底変換写像を与える
（正確な主張については証明を参照されたい）。
\end{lemma}

\begin{proof}
仮定から $g \circ f$ と $m$ は Tor 独立である
（詳細は省略する）ので、この主張には意味がある。
この証明では $Lk^*$ の代わりに $k^*$、$Rf_*$ の代わりに
$f_*$ と書く。補題 \ref{spaces-duality-lemma-twisted-inverse-image} による
$f$, $g$, $g \circ f$ の右随伴を $a$, $b$, $c$ とし、
プライム付きの場合も同様に書く。上の正方形に対応する矢印は合成
$$
\gamma_{top} :
k^* \circ a \to k^* \circ a \circ l_* \circ l^*
\xleftarrow{\xi_{top}} k^* \circ k_* \circ a' \circ l^* \to a' \circ l^*
$$
である。ここで $\xi_{top} : k_* \circ a' \to a \circ l_*$ は
同型（したがって逆にできる）であり、基底変換写像
$l^* \circ f_* \to f'_* \circ k^*$ に「双対」な矢印である。
外側の矢印は標準写像 $1 \to l_* \circ l^*$ および
$k^* \circ k_* \to 1$ から来る。同様に、下の正方形については
$$
\gamma_{bot} :
l^* \circ b \to l^* \circ b \circ m_* \circ m^*
\xleftarrow{\xi_{bot}} l^* \circ l_* \circ b' \circ m^* \to b' \circ m^*
$$
を得る。外側の長方形については
$$
\gamma_{rect} :
k^* \circ c \to k^* \circ c \circ m_* \circ m^*
\xleftarrow{\xi_{rect}} k^* \circ k_* \circ c' \circ m^* \to c' \circ m^*
$$
を得る。$(g \circ f)_* = g_* \circ f_*$ であるから
$c = a \circ b$ であり、同様に $c' = a' \circ b'$ である。
補題の主張は、$\gamma_{rect}$ が合成
$$
k^* \circ c = k^* \circ a \circ b \xrightarrow{\gamma_{top}}
a' \circ l^* \circ b \xrightarrow{\gamma_{bot}}
a' \circ b' \circ m^* = c' \circ m^*
$$
に等しいということである。これを見るため、次の図式を考える：
$$
\xymatrix@C=.8em{
& & k^* \circ a \circ b \ar[d] \ar[lldd] \\
& & k^* \circ a \circ l_* \circ l^* \circ b \ar[ld] \\
k^* \circ a \circ b \circ m_* \circ m^* \ar[r] &
k^* \circ a \circ l_* \circ l^* \circ b \circ m_* \circ m^* &
k^* \circ k_* \circ a' \circ l^* \circ b \ar[u]_{\xi_{top}} \ar[d] \ar[ld] \\
& k^*\circ k_* \circ a' \circ l^* \circ b \circ m_* \circ m^*
\ar[u]_{\xi_{top}} \ar[rd] &
a' \circ l^* \circ b \ar[d] \\
k^* \circ k_* \circ a' \circ b' \circ m^* \ar[uu]_{\xi_{rect}} \ar[ddrr] &
k^*\circ k_* \circ a' \circ l^* \circ l_* \circ b' \circ m^*
\ar[u]_{\xi_{bot}} \ar[l] \ar[dr] &
a' \circ l^* \circ b \circ m_* \circ m^* \\
& & a' \circ l^* \circ l_* \circ b' \circ m^* \ar[u]_{\xi_{bot}} \ar[d] \\
& & a' \circ b' \circ m^*
}
$$
右辺を下ると上記の合成が得られ、左辺を下ると
$\gamma_{rect}$ が得られる。この図式の右辺にあるすべての四辺形は、
『圏』補題 \ref{categories-lemma-properties-2-cat-cats}、またはより直接には
『圏』定義 \ref{categories-definition-horizontal-composition}
に先立つ議論により可換である。したがって、図式
$$
\xymatrix{
a \circ l_* \circ l^* \circ b \circ m_* &
a \circ b \circ m_* \ar[l] \\
k_* \circ a' \circ l^* \circ b \circ m_* \ar[u]_{\xi_{top}} & \\
k_* \circ a' \circ l^* \circ l_* \circ b' \ar[u]_{\xi_{bot}} \ar[r] &
k_* \circ a' \circ b' \ar[uu]_{\xi_{rect}}
}
$$
が、矢印 $\xi_{top}$, $\xi_{bot}$, $\xi_{rect}$ を逆にすると
可換になることを示せば十分である
（これは、この図式そのものが可換であると要求することとは異なる）。
ところが図式
$$
\xymatrix{
& a \circ l_* \circ l^* \circ b \circ m_* \\
a \circ l_* \circ l^* \circ l_* \circ b'
\ar[ru]^{\xi_{bot}} & &
k_* \circ a' \circ l^* \circ b \circ m_* \ar[ul]_{\xi_{top}} \\
& k_* \circ a' \circ l^* \circ l_* \circ b'
\ar[ul]^{\xi_{top}} \ar[ur]_{\xi_{bot}}
}
$$
は『圏』補題 \ref{categories-lemma-properties-2-cat-cats} により可換である。
次の二つの図式も可換であり（上記の参考文献を参照）、
% P10-E0257
$$
\vcenter{
\xymatrix@C=.8em{
a \circ l_* \circ l^* \circ b \circ m_* & a \circ b \circ m \ar[l] \\
a \circ l_* \circ l^* \circ l_* \circ b' \ar[u] &
a \circ l_* \circ b' \ar[l] \ar[u]
}
}
\quad\text{および}\quad
\vcenter{
\xymatrix@C=.8em{
a \circ l_* \circ l^* \circ l_* \circ b' \ar[r] & a \circ l_* \circ b' \\
k_* \circ a' \circ l^* \circ l_* \circ b' \ar[u] \ar[r] &
k_* \circ a' \circ b' \ar[u]
}
}
$$
$l_* \to l_* \circ l^* \circ l_* \to l_*$ の合成は恒等写像なので、
% P10-E0258
$$
k \circ a' \circ b' \xrightarrow{\xi_{bot}} a \circ l_* \circ b
\xrightarrow{\xi_{top}} a \circ b \circ m_*
$$
が $\xi_{rect}$ に等しいことを証明すれば十分である
（同一視 $a \circ b = c$ および $a' \circ b' = c'$ を介して）。
これは『サイト上のコホモロジー』評注
\ref{sites-cohomology-remark-compose-base-change} に双対な主張であり、
証明は完了する。
\end{proof}

\begin{lemma}
\label{spaces-duality-lemma-compose-base-change-maps-horizontal}
$S$ をスキームとする。$S$ 上の準コンパクトかつ準分離な代数空間の
可換図式
$$
\xymatrix{
X'' \ar[r]_{g'} \ar[d]_{f''} & X' \ar[r]_g \ar[d]_{f'} & X \ar[d]^f \\
Y'' \ar[r]^{h'} & Y' \ar[r]^h & Y
}
$$
を考える。両方の正方形は Cartesian であり、$f$ と $h$、
および $f'$ と $h'$ は Tor 独立であるとする。このとき、
二つの正方形に対する写像 (\ref{spaces-duality-equation-base-change-map}) の合成は、
外側の長方形に対する基底変換写像を与える
（正確な主張については証明を参照されたい）。
\end{lemma}

\begin{proof}
仮定から $f$ と $h \circ h'$ は Tor 独立である
（詳細は省略する）ので、この主張には意味がある。
この証明では $Lg^*$ の代わりに $g^*$、$Rf_*$ の代わりに
$f_*$ と書く。補題 \ref{spaces-duality-lemma-twisted-inverse-image} による
$f$, $f'$, $f''$ の右随伴を $a$, $a'$, $a''$ とする。
右の正方形に対応する矢印は合成
$$
\gamma_{right} :
g^* \circ a \to g^* \circ a \circ h_* \circ h^*
\xleftarrow{\xi_{right}} g^* \circ g_* \circ a' \circ h^* \to a' \circ h^*
$$
である。ここで $\xi_{right} : g_* \circ a' \to a \circ h_*$ は
同型（したがって逆にできる）であり、基底変換写像
$h^* \circ f_* \to f'_* \circ g^*$ に「双対」な矢印である。
外側の矢印は標準写像 $1 \to h_* \circ h^*$ および
$g^* \circ g_* \to 1$ から来る。同様に、左の正方形については
$$
\gamma_{left} :
(g')^* \circ a' \to (g')^* \circ a' \circ (h')_* \circ (h')^*
\xleftarrow{\xi_{left}}
(g')^* \circ (g')_* \circ a'' \circ (h')^* \to a'' \circ (h')^*
$$
を得る。外側の長方形については
$$
\gamma_{rect} :
k^* \circ a \to
k^* \circ a \circ m_* \circ m^* \xleftarrow{\xi_{rect}}
k^* \circ k_* \circ a'' \circ m^* \to
a'' \circ m^*
$$
を得る。ここで $k = g \circ g'$ および $m = h \circ h'$ である。
$k^* = (g')^* \circ g^*$ および $m^* = (h')^* \circ h^*$ である。
補題の主張は、$\gamma_{rect}$ が合成
$$
k^* \circ a =
(g')^* \circ g^* \circ a \xrightarrow{\gamma_{right}}
(g')^* \circ a' \circ h^* \xrightarrow{\gamma_{left}}
a'' \circ (h')^* \circ h^* = a'' \circ m^*
$$
に等しいということである。これを見るため、次の図式を考える：
$$
\xymatrix{
& (g')^* \circ g^* \circ a \ar[d] \ar[ddl] \\
& (g')^* \circ g^* \circ a \circ h_* \circ h^* \ar[ld] \\
(g')^* \circ g^* \circ a \circ h_* \circ (h')_* \circ (h')^* \circ h^* &
(g')^* \circ g^* \circ g_* \circ a' \circ h^*
\ar[u]_{\xi_{right}} \ar[d] \ar[ld] \\
(g')^* \circ g^* \circ g_* \circ a' \circ (h')_* \circ (h')^* \circ h^*
\ar[u]_{\xi_{right}} \ar[dr] &
(g')^* \circ a' \circ h^* \ar[d] \\
(g')^* \circ g^* \circ g_* \circ (g')_* \circ a'' \circ (h')^* \circ h^*
\ar[u]_{\xi_{left}} \ar[ddr] \ar[dr] &
(g')^* \circ a' \circ (h')_* \circ (h')^* \circ h^* \\
& (g')^*\circ (g')_* \circ a'' \circ (h')^* \circ h^*
\ar[u]_{\xi_{left}} \ar[d] \\
& a'' \circ (h')^* \circ h^*
}
$$
右辺を下ると上記の合成が得られ、左辺を下ると
$\gamma_{rect}$ が得られる。この図式の右辺にあるすべての四辺形は、
『圏』補題 \ref{categories-lemma-properties-2-cat-cats}、またはより直接には
『圏』定義 \ref{categories-definition-horizontal-composition}
に先立つ議論により可換である。したがって
$$
g_* \circ (g')_* \circ a'' \xrightarrow{\xi_{left}}
g_* \circ a' \circ (h')_* \xrightarrow{\xi_{right}}
a \circ h_* \circ (h')_*
$$
が $\xi_{rect}$ に等しいことを示せば十分である。
これは『コホモロジー』評注
\ref{cohomology-remark-compose-base-change-horizontal} に双対な主張であり、
証明は完了する。
\end{proof}

\begin{remark}
\label{spaces-duality-remark-going-around}
$S$ をスキームとする。$S$ 上の準コンパクトかつ準分離な代数空間の
可換図式
$$
\xymatrix{
X'' \ar[r]_{k'} \ar[d]_{f''} & X' \ar[r]_k \ar[d]_{f'} & X \ar[d]^f \\
Y'' \ar[r]^{l'} \ar[d]_{g''} & Y' \ar[r]^l \ar[d]_{g'} & Y \ar[d]^g \\
Z'' \ar[r]^{m'} & Z' \ar[r]^m & Z
}
$$
を考える。すべての正方形は Cartesian であり、
$(f, l)$, $(g, m)$, $(f', l')$, $(g', m')$ は
Tor 独立な射の組であるとする。補題
\ref{spaces-duality-lemma-twisted-inverse-image} による
$f$, $f'$, $f''$, $g$, $g'$, $g''$ の右随伴をそれぞれ
$a$, $a'$, $a''$, $b$, $b'$, $b''$ とする。
図式の正方形を次のように $A$, $B$, $C$, $D$ と名付ける：
$$
\begin{matrix}
A & B \\
C & D
\end{matrix}
$$
各正方形に対する写像 (\ref{spaces-duality-equation-base-change-map}) は
（ここでは $k^* = Lk^*$ などと書く）
$$
\begin{matrix}
\gamma_A : (k')^* \circ a' \to a'' \circ (l')^* &
\gamma_B : k^* \circ a \to a' \circ l^* \\
\gamma_C : (l')^* \circ b' \to b'' \circ (m')^* &
\gamma_D : l^* \circ b \to b' \circ m^*
\end{matrix}
$$
である。$2 \times 1$ および $1 \times 2$ の長方形に対しては、
さらに四つの基底変換写像
$$
\begin{matrix}
\gamma_{A + B} : (k \circ k')^* \circ a \to a'' \circ (l \circ l')^* \\
\gamma_{C + D} : (l \circ l')^* \circ b \to b'' \circ (m \circ m')^* \\
\gamma_{A + C} : (k')^* \circ (a' \circ b') \to (a'' \circ b'') \circ (m')^* \\
% P10-E0259
\gamma_{A + C} : k^* \circ (a \circ b) \to (a' \circ b') \circ m^*
\end{matrix}
$$
がある。補題 \ref{spaces-duality-lemma-compose-base-change-maps-horizontal} により
$$
\gamma_{A + B} = \gamma_A \circ \gamma_B, \quad
\gamma_{C + D} = \gamma_C \circ \gamma_D
$$
であり、補題 \ref{spaces-duality-lemma-compose-base-change-maps} により
$$
\gamma_{A + C} = \gamma_C \circ \gamma_A, \quad
\gamma_{B + D} = \gamma_D \circ \gamma_B
$$
である。ここで、より正確には、『圏』節
\ref{categories-section-formal-cat-cat} の記法を用いて
$\gamma_{A + B} = (\gamma_A \star \text{id}_{l^*}) \circ
(\text{id}_{(k')^*} \star \gamma_B)$ と書くべきであり、
他についても同様である。しかし、補題
\ref{spaces-duality-lemma-compose-base-change-maps} および
\ref{spaces-duality-lemma-compose-base-change-maps-horizontal} の証明で用いた記法の濫用、
すなわち恒等変換との $\star$ 積を省略する慣行を続ける。
変換の始域と終域が分かっていれば、どれを補うべきか判別できるからである。
% P10-E0260
以上を踏まえると、（先験的には）二つの変換
$$
(k')^* \circ k^* \circ a \circ b
\longrightarrow
a'' \circ b'' \circ (m')^* \circ m^*
$$
を得る。すなわち
$$
\gamma_C \circ \gamma_A \circ \gamma_D \circ \gamma_B =
\gamma_{A + C} \circ \gamma_{B + D}
$$
および
$$
\gamma_C \circ \gamma_D \circ \gamma_A \circ \gamma_B =
\gamma_{C + D} \circ \gamma_{A + B}
$$
である。この評注の要点は、これらの変換が等しいことを指摘することにある。
実際、そのためには
$$
\xymatrix{
(k')^* \circ a' \circ l^* \circ b \ar[r]_{\gamma_D} \ar[d]_{\gamma_A} &
(k')^* \circ a' \circ b' \circ m^* \ar[d]^{\gamma_A} \\
a'' \circ (l')^* \circ l^* \circ b \ar[r]^{\gamma_D} &
a'' \circ (l')^* \circ b' \circ m^*
}
$$
が可換であることを示せば十分である。これは『圏』補題
\ref{categories-lemma-properties-2-cat-cats}、またはより直接には
『圏』定義 \ref{categories-definition-horizontal-composition}
に先立つ議論により成り立つ。
\end{remark}






\section{順像の右随伴と基底変換 II}
\label{spaces-duality-section-base-change-II}

\noindent
本節では、節 \ref{spaces-duality-section-base-change-map} の基底変換写像が
いくつかの場合に同型であることを証明する。

\begin{lemma}
\label{spaces-duality-lemma-more-base-change}
図式 (\ref{spaces-duality-equation-base-change}) において、さらに
$g : Y' \to Y$ がアフィン・スキームの射であり、
$f : X \to Y$ が固有であると仮定する。このとき基底変換写像
(\ref{spaces-duality-equation-base-change-map}) は、次の場合に同型
$$
L(g')^*a(K) \longrightarrow a'(Lg^*K)
$$
を誘導する。
\begin{enumerate}
% P10-E0261
\item $f$ が有限表示平坦なら、すべての
$K \in D_\QCoh(\mathcal{O}_X)$ に対して。
\item $f$ が完全で $Y$ が Noether なら、すべての
$K \in D_\QCoh(\mathcal{O}_X)$ に対して。
\item $g$ の Tor 次元が有限で $Y$ が Noether なら、
$K \in D_\QCoh^+(\mathcal{O}_X)$ に対して。
\end{enumerate}
\end{lemma}

\begin{proof}
$Y = \Spec(A)$ および $Y' = \Spec(A')$ と書く。
アフィン射の基底変換なので、射 $g'$ はアフィンである。
$M$ を $D_\QCoh(\mathcal{O}_X)$ の完全生成対象とする。
『空間の導来圏』定理
\ref{spaces-perfect-theorem-bondal-van-den-Bergh} を参照されたい。
このとき $L(g')^*M$ は $D_\QCoh(\mathcal{O}_{X'})$ の生成対象である。
『空間の導来圏』評注
\ref{spaces-perfect-remark-pullback-generator} を参照されたい。
したがって、(\ref{spaces-duality-equation-base-change-map}) が大域 Hom 複体の同型
\begin{equation}
\label{spaces-duality-equation-iso}
R\Hom_{X'}(L(g')^*M, L(g')^*a(K))
\longrightarrow
R\Hom_{X'}(L(g')^*M, a'(Lg^*K))
\end{equation}
を誘導することを示せば十分である。『サイト上のコホモロジー』節
\ref{sites-cohomology-section-global-RHom} を参照されたい。
実際、これにより $L(g')^*a(K) \to a'(Lg^*K)$ の錐が零であることが従う。
証明の構成は次のとおりである。まずこれらの Hom 複体が同型であることを示し、
証明の最後の部分で、その同型が (\ref{spaces-duality-equation-iso}) により誘導されることを示す。

\medskip\noindent
左辺について考える。$M$ は完全なので、標準写像
$$
R\Hom_X(M, a(K)) \otimes^\mathbf{L}_A A'
\longrightarrow
R\Hom_{X'}(L(g')^*M, L(g')^*a(K))
$$
は『空間の導来圏』補題
\ref{spaces-perfect-lemma-affine-morphism-and-hom-out-of-perfect}
により同型である。これと補題 \ref{spaces-duality-lemma-iso-global-hom} の同型
$R\Hom_Y(Rf_*M, K) = R\Hom_X(M, a(K))$ を組み合わせると、
左辺は $R\Hom_Y(Rf_*M, K) \otimes^\mathbf{L}_A A'$ に等しい。

\medskip\noindent
右辺について考える。まず補題 \ref{spaces-duality-lemma-iso-global-hom} の同型
$$
R\Hom_{X'}(L(g')^*M, a'(Lg^*K)) = R\Hom_{Y'}(Rf'_*L(g')^*M, Lg^*K)
$$
を用いる。$f$ と $g$ は Tor 独立なので、次の基底変換写像を考える。
\par\noindent
$Lg^*Rf_*M \to Rf'_*L(g')^*M$。
これは『空間の導来圏』補題
\ref{spaces-perfect-lemma-compare-base-change} により同型である。
したがって、これを $R\Hom_{Y'}(Lg^*Rf_*M,\allowbreak Lg^*K)$ と書き換えられる。
$Y$, $Y'$ はアフィンであり、$K$, $Rf_*M$ は
$D_\QCoh(\mathcal{O}_Y)$ に属する
（『空間の導来圏』補題
\ref{spaces-perfect-lemma-quasi-coherence-direct-image}）ので、
$D(A')$ における標準写像
$$
\beta :
R\Hom_Y(Rf_*M, K) \otimes^\mathbf{L}_A A'
\longrightarrow
R\Hom_{Y'}(Lg^*Rf_*M, Lg^*K)
$$
がある。これは『代数詳論』式
(\ref{more-algebra-equation-base-change-RHom}) の矢印である。
ここでは『スキームの導来圏』補題
\ref{perfect-lemma-affine-compare-bounded} および
\ref{perfect-lemma-quasi-coherence-internal-hom} を用いて、
代数との間を行き来した。
\begin{enumerate}
\item $f$ が平坦かつ有限表示なら、『空間の導来圏』補題
\ref{spaces-perfect-lemma-flat-proper-perfect-direct-image-general}
により複体 $Rf_*M$ は $Y$ 上完全であり、『代数詳論』補題
\ref{more-algebra-lemma-base-change-RHom} の (1) により
$\beta$ は同型である。
\item $f$ が完全で $Y$ が Noether なら、『空間の射詳論』補題
\ref{spaces-more-morphisms-lemma-perfect-proper-perfect-direct-image}
により複体 $Rf_*M$ は $Y$ 上完全であり、先ほどと同様に
$\beta$ は同型である。
\item $g$ の Tor 次元が有限で $Y$ が Noether なら、
『空間の導来圏』補題
\ref{spaces-perfect-lemma-direct-image-coherent} および
\ref{spaces-perfect-lemma-identify-pseudo-coherent-noetherian}
により複体 $Rf_*M$ は $Y$ 上擬連接であり、『代数詳論』補題
\ref{more-algebra-lemma-base-change-RHom} の (4) により
$\beta$ は同型である。
\end{enumerate}
前段落と同じ答えが得られた。

\medskip\noindent
証明の残りでは、第二、第三段落で与えた
(\ref{spaces-duality-equation-iso}) の左辺と右辺の同一視が、実際に
(\ref{spaces-duality-equation-iso}) により与えられることを示す。
式を扱いやすくするため、$(-, -)_X = R\Hom_X(-, -)$、
$- \otimes_A^\mathbf{L} A'$ の代わりに $- \otimes A'$、
% P10-E0262
$g^* = Lg^*$ および $f_* = Rf_*$ と略記する。次の可換図式を考える：
$$
\xymatrix@C=.4em{
((g')^*M, (g')^*a(K))_{X'} \ar[d] &
(M, a(K))_X \otimes A' \ar[l]^-\alpha \ar[d] &
(f_*M, K)_Y \otimes A' \ar@{=}[l] \ar[d] \\
((g')^*M, (g')^*a(g_*g^*K))_{X'} &
(M, a(g_*g^*K))_X \otimes A' \ar[l]^-\alpha &
(f_*M, g_*g^*K)_Y \otimes A' \ar@{=}[l] \ar@/_4pc/[dd]_{\mu'} \\
((g')^*M, (g')^*g'_*a'(g^*K))_{X'} \ar[u] \ar[d] &
(M, g'_*a'(g^*K))_X \otimes A' \ar[u] \ar[l]^-\alpha \ar[ld]^\mu &
% P10-E0263
(f_*M, K) \otimes A' \ar[d]^\beta \\
((g')^*M, a'(g^*K))_{X'} &
(f'_*(g')^*M, g^*K)_{Y'} \ar@{=}[l] \ar[r] &
(g^*f_*M, g^*K)_{Y'}
}
$$
$\alpha$ と記した矢印は、角が $X', X, Y', Y$ である図式に対する
『空間の導来圏』補題
\ref{spaces-perfect-lemma-affine-morphism-and-hom-out-of-perfect}
の写像である。水平矢印は各変数について関手的なので、図式の上部は可換である。
中央の垂直矢印は、補題 \ref{spaces-duality-lemma-flat-precompose-pus} の可逆変換
$g'_* \circ a' \to a \circ g_*$ から来る。したがって中央の正方形は可換である。
左辺を下る写像が (\ref{spaces-duality-equation-iso}) である。
上の水平矢印は、証明の第二段落で用いた同一視を与える。
$\beta$ を含む下の水平矢印は、証明の第三段落で用いた同一視を与える。
$E \in D(A)$, $E' \in D(A')$ および $D(A)$ における
$c : E \to E'$ が与えられたとき、忘却と基底変換の随伴性により
$c$ から誘導される写像を $\mu_c : E \otimes A' \to E'$ と書く。
$c$ が明らかなら $\mu = \mu_c$ と書き、
記法から $c$ を省く。図式中の $\mu$ は、この形の写像であり、
% P10-E0264
同一視 $(M, g'_*a(g^*K))_X = ((g')^*M, a'(g^*K))_{X'}$
によって $c$ が与えられる。$\mu$ を含む三角形は『空間の導来圏』評注
\ref{spaces-perfect-remark-multiplication-map} により可換である。

\medskip\noindent
次の図式が可換であることに注意する：
$$
\xymatrix{
(M, a(g_*g^*K))_X &
(f_*M, g_* g^*K)_Y \ar@{=}[l] &
(g^*f_*M, g^*K)_{Y'} \ar@{=}[l] \\
(M, g'_* a'(g^*K))_X \ar[u] &
((g')^*M, a'(g^*K))_{X'} \ar@{=}[l] &
(f'_*(g')^*M, g^*K)_{Y'} \ar@{=}[l] \ar[u]
}
$$
これは変換 $g'_* \circ a' \to a \circ g_*$ の定義そのものによる。
上記と同様に、同一視
% P10-E0265
$(f_*M, g_*g^*K)_X = (g^*f_*M, g^*K)_{Y'}$ に対応する
写像を $\mu'$ とすると、六角形も可換である。したがって、
$\beta$ が
% P10-E0266
$(f_*M, K)_Y \otimes A' \to (f_*M, g_*g^*K)_X \otimes A'$
と $\mu'$ の合成に等しいことを示せば十分である。
そのためには、誘導される二つの写像
$(f_*M, K)_Y \to (g^*f_*M, g^*K)_{Y'}$ が同じであることを示せばよい。
言い換えると、すべての $E, K \in D(A)$ に対して図式
$$
\xymatrix@C=.8em{
R\Hom_A(E, K) \ar[rr]_{\beta\text{ により誘導}} \ar[rd] & &
R\Hom_{A'}(E \otimes_A^\mathbf{L} A', K \otimes_A^\mathbf{L} A') \\
& R\Hom_A(E, K \otimes_A^\mathbf{L} A') \ar[ru]
}
$$
が可換であることを示せば十分である。これは『代数詳論』節
\ref{more-algebra-section-base-change-RHom} における $\beta$ の構成そのもの
なので、証明は完了する。
\end{proof}






\section{順像の右随伴とトレース写像}
\label{spaces-duality-section-trace}

\noindent
$S$ をスキームとする。
$f : X \to Y$ を、$S$ 上の準コンパクトかつ準分離な代数空間の射とする。
$a : D_\QCoh(\mathcal{O}_Y) \to D_\QCoh(\mathcal{O}_X)$ を、
補題 \ref{spaces-duality-lemma-twisted-inverse-image} における右随伴とする。
『圏』節 \ref{categories-section-adjoint} により、関手の変換
$$
\text{Tr}_f : Rf_* \circ a \longrightarrow \text{id}
$$
を得る。$K \in D_\QCoh(\mathcal{O}_Y)$ に対する対応する写像
$\text{Tr}_{f, K} : Rf_*a(K) \longrightarrow K$ を、しばしば
{\it トレース写像}という。この写像は、右随伴を特徴付ける全単射
$$
\Hom_X(L, a(K)) \longrightarrow \Hom_Y(Rf_*L, K)
$$
が、$L \in D_\QCoh(\mathcal{O}_X)$ に対して
$$
\varphi \longmapsto \text{Tr}_{f, K} \circ Rf_*\varphi
$$
で与えられるという性質をもつ。標準写像
(\ref{spaces-duality-equation-sheafy-trace})
$$
Rf_*R\SheafHom_{\mathcal{O}_X}(L, a(K))
\longrightarrow
R\SheafHom_{\mathcal{O}_Y}(Rf_*L, K)
$$
は $\text{Tr}_{f, K}$ との合成により生じる。
本節で考えるすべてのトレース写像は、このトレース写像の特殊な場合である。
いくつかの特殊な場合を論じる前に、トレース写像の形成が基底変換と可換する
ことを示す。

\begin{lemma}[トレース写像と基底変換]
\label{spaces-duality-lemma-trace-map-and-base-change}
図式 (\ref{spaces-duality-equation-base-change}) が与えられているとする。このとき、写像
$1 \star \text{Tr}_f : Lg^* \circ Rf_* \circ a \to Lg^*$ および
$\text{Tr}_{f'} \star 1 : Rf'_* \circ a' \circ Lg^* \to Lg^*$ は、
基底変換写像
$\beta : Lg^* \circ Rf_* \to Rf'_* \circ L(g')^*$
（『サイト上のコホモロジー』評注
\ref{sites-cohomology-remark-base-change}）および
$\alpha : L(g')^* \circ a \to a' \circ Lg^*$
（(\ref{spaces-duality-equation-base-change-map})）を介して一致する。
より正確には、関手の変換の図式
$$
\xymatrix{
Lg^* \circ Rf_* \circ a
\ar[d]_{\beta \star 1} \ar[r]_-{1 \star \text{Tr}_f} &
Lg^* \\
Rf'_* \circ L(g')^* \circ a \ar[r]^{1 \star \alpha} &
Rf'_* \circ a' \circ Lg^* \ar[u]_{\text{Tr}_{f'} \star 1}
}
$$
は可換である。
\end{lemma}

\begin{proof}
この証明では $Rf_*$ の代わりに $f_*$、$Lg^*$ の代わりに
$g^*$ と書き、恒等変換との $\star$ 積を省略する。% P10-E0268
変換の始域と終域が分かっていれば、どれを補うべきか判別できるからである。
$\beta : g^* \circ f_* \to f'_* \circ (g')^*$ は同型であり、
$\alpha$ は $\beta$ の随伴である同型
$\beta^\vee : g'_* \circ a' \to a \circ g_*$ を用いて定義されることを
思い出そう。補題 \ref{spaces-duality-lemma-flat-precompose-pus} とその証明を参照されたい。
まず、補題の図式の上の水平矢印は合成
$$
g^* \circ f_* \circ a \to
g^* \circ f_* \circ a \circ g_* \circ g^* \to
g^* \circ g_* \circ g^* \to g^*
$$
に等しいことに注意する。ここで最初の矢印は $(g^*, g_*)$ の単位、
二番目の矢印は $\text{Tr}_f$、三番目の矢印は $(g^*, g_*)$ の余単位である。
これは、単位と余単位の合成
$g^* \to g^* \circ g_* \circ g^* \to g^*$ が恒等写像であることの
簡単な帰結である。図式
$$
\xymatrix@C=.3em{
& g^* \circ f_* \circ a \ar[ld]_\beta \ar[d] \ar[r]_{\text{Tr}_f} & g^* \\
f'_* \circ (g')^* \circ a \ar[dr] &
g^* \circ f_* \circ a \circ g_* \circ g^* \ar[d]_\beta \ar[ru] &
g^* \circ f_* \circ g'_* \circ a' \circ g^* \ar[l]_{\beta^\vee} \ar[d]_\beta &
f'_* \circ a' \circ g^* \ar[lu]_{\text{Tr}_{f'}} \\
& f'_* \circ (g')^* \circ a \circ g_* \circ g^* &
f'_* \circ (g')^* \circ g'_* \circ a' \circ g^* \ar[ru] \ar[l]_{\beta^\vee}
}
$$
を考える。この図式の二つの正方形は、『圏』補題
\ref{categories-lemma-properties-2-cat-cats}、またはより直接には
『圏』定義 \ref{categories-definition-horizontal-composition}
に先立つ議論により可換である。% P10-E0267
三角形は上の議論により可換である。『圏』補題
\ref{categories-lemma-transformation-between-functors-and-adjoints}
により、正方形
$$
\xymatrix{
g^* \circ f_* \circ g'_* \circ a' \ar[d]_{\beta^\vee} \ar[r]_-\beta &
f'_* \circ (g')^* \circ g'_* \circ a' \ar[d] \\
g^* \circ f_* \circ a \circ g_* \ar[r] &
\text{id}
}
$$
は可換であり、これは大きな図式の五角形が可換であることを含意する。
$\beta$ と $\beta^\vee$ は同型であり、また大きな図式の外周を進む合成は、
% P10-E0269
定義により $\text{Tr}_f \circ \alpha \circ \beta$ に等しいので、
補題が証明された。
\end{proof}

\noindent
$S$ をスキームとする。
$f : X \to Y$ を、$S$ 上の準コンパクトかつ準分離な代数空間の射とする。
$a : D_\QCoh(\mathcal{O}_Y) \to D_\QCoh(\mathcal{O}_X)$ を、
補題 \ref{spaces-duality-lemma-twisted-inverse-image} における $Rf_*$ の右随伴とする。
『圏』節 \ref{categories-section-adjoint} により、随伴の単位と呼ばれる
関手の変換
$$
\eta_f : \text{id} \to  a \circ Rf_*
$$
を得る。

\begin{lemma}
\label{spaces-duality-lemma-unit-and-base-change}
図式 (\ref{spaces-duality-equation-base-change}) が与えられているとする。このとき、写像
$1 \star \eta_f : L(g')^* \to L(g')^* \circ a \circ Rf_*$ および
$\eta_{f'} \star 1 : L(g')^* \to a' \circ Rf'_* \circ L(g')^*$ は、
基底変換写像
$\beta : Lg^* \circ Rf_* \to Rf'_* \circ L(g')^*$
（『サイト上のコホモロジー』評注
\ref{sites-cohomology-remark-base-change}）および
$\alpha : L(g')^* \circ a \to a' \circ Lg^*$
（(\ref{spaces-duality-equation-base-change-map})）を介して一致する。
より正確には、関手の変換の図式
$$
\xymatrix{
L(g')^* \ar[r]_-{1 \star \eta_f} \ar[d]_{\eta_{f'} \star 1} &
L(g')^* \circ a \circ Rf_* \ar[d]^\alpha \\
a' \circ Rf'_* \circ L(g')^* &
a' \circ Lg^* \circ Rf_* \ar[l]_-\beta
}
$$
は可換である。
\end{lemma}

\begin{proof}
この証明は補題 \ref{spaces-duality-lemma-trace-map-and-base-change} の証明に双対である。
この証明では $Rf_*$ の代わりに $f_*$、$Lg^*$ の代わりに
$g^*$ と書き、恒等変換との $\star$ 積を省略する。% P10-E0268
変換の始域と終域が分かっていれば、どれを補うべきか判別できるからである。
$\beta : g^* \circ f_* \to f'_* \circ (g')^*$ は同型であり、
$\alpha$ は $\beta$ の随伴である同型
$\beta^\vee : g'_* \circ a' \to a \circ g_*$ を用いて定義されることを
思い出そう。補題 \ref{spaces-duality-lemma-flat-precompose-pus} とその証明を参照されたい。
まず、補題の図式の左の垂直矢印は合成
$$
(g')^* \to (g')^* \circ g'_* \circ (g')^* \to
(g')^* \circ g'_* \circ a' \circ f'_* \circ (g')^* \to
a' \circ f'_* \circ (g')^*
$$
に等しいことに注意する。ここで最初の矢印は $((g')^*, g'_*)$ の単位、
二番目の矢印は $\eta_{f'}$、三番目の矢印は
$((g')^*, g'_*)$ の余単位である。これは、単位と余単位の合成
$(g')^* \to (g')^* \circ (g')_* \circ (g')^* \to (g')^*$ が
恒等写像であることの簡単な帰結である。図式
$$
\xymatrix{
& (g')^* \circ a \circ f_* \ar[r] &
(g')^* \circ a \circ g_* \circ g^* \circ f_*
\ar[ld]_\beta \\
(g')^* \ar[ru]^{\eta_f} \ar[dd]_{\eta_{f'}} \ar[rd] &
(g')^* \circ a \circ g_* \circ f'_* \circ (g')^* &
(g')^* \circ g'_* \circ a' \circ g^* \circ f_*
\ar[u]_{\beta^\vee} \ar[ld]_\beta \ar[d] \\
& (g')^* \circ g'_* \circ a' \circ f'_* \circ (g')^*
\ar[ld] \ar[u]_{\beta^\vee} &
a' \circ g^* \circ f_* \ar[lld]^\beta \\
a' \circ f'_* \circ (g')^*
}
$$
を考える。この図式の二つの正方形は、『圏』補題
\ref{categories-lemma-properties-2-cat-cats}、またはより直接には
『圏』定義 \ref{categories-definition-horizontal-composition}
に先立つ議論により可換である。% P10-E0267
三角形は上の議論により可換である。『圏』補題
\ref{categories-lemma-transformation-between-functors-and-adjoints}
の双対により、正方形
$$
\xymatrix{
\text{id} \ar[r] \ar[d] &
g'_* \circ a' \circ g^* \circ f_* \ar[d]^\beta \\
g'_* \circ a' \circ g^* \circ f_* \ar[r]^{\beta^\vee} &
a \circ g_* \circ f'_* \circ (g')^*
}
$$
は可換であり、これは大きな図式の五角形が可換であることを含意する。
$\beta$ と $\beta^\vee$ は同型であり、また大きな図式の外周を進む合成は、
定義により $\beta \circ \alpha \circ \eta_f$ に等しいので、補題が証明された。
\end{proof}






\section{順像の右随伴と引戻し}
\label{spaces-duality-section-compare-with-pullback}

\noindent
$S$ をスキームとする。
$f : X \to Y$ を、$S$ 上の準コンパクトかつ準分離な代数空間の射とする。
$a$ を補題 \ref{spaces-duality-lemma-twisted-inverse-image} における順像の右随伴とする。
$K, L \in D_\QCoh(\mathcal{O}_Y)$ に対して標準写像
$$
Lf^*K \otimes^\mathbf{L}_{\mathcal{O}_X} a(L)
\longrightarrow
a(K \otimes_{\mathcal{O}_Y}^\mathbf{L} L)
$$
がある。すなわち、この写像は写像
$$
Rf_*(Lf^*K \otimes^\mathbf{L}_{\mathcal{O}_X} a(L)) =
K \otimes^\mathbf{L}_{\mathcal{O}_Y} Rf_*(a(L))
\longrightarrow
K \otimes^\mathbf{L}_{\mathcal{O}_Y} L
$$
の随伴である（等号は『空間の導来圏』補題
\ref{spaces-perfect-lemma-cohomology-base-change} による）。
ここではトレース写像 $Rf_*a(L) \to L$ を用いる。
$L = \mathcal{O}_Y$ のとき、$K$ について関手的で、識別三角形と両立する写像
\begin{equation}
\label{spaces-duality-equation-compare-with-pullback}
Lf^*K \otimes^\mathbf{L}_{\mathcal{O}_X} a(\mathcal{O}_Y) \longrightarrow a(K)
\end{equation}
を得る。

\begin{lemma}
\label{spaces-duality-lemma-compare-with-pullback-perfect}
$S$ をスキームとする。
$f : X \to Y$ を、$S$ 上の準コンパクトかつ準分離な代数空間の射とする。
上で $K, L \in D_\QCoh(\mathcal{O}_Y)$ に対して定義した写像
$Lf^*K \otimes^\mathbf{L}_{\mathcal{O}_X} a(L) \to
a(K \otimes_{\mathcal{O}_Y}^\mathbf{L} L)$ は、$K$ が完全なら同型である。
特に、$K$ が完全なら (\ref{spaces-duality-equation-compare-with-pullback}) は同型である。
\end{lemma}

\begin{proof}
$K^\vee$ を $K$ の「双対」とする。『サイト上のコホモロジー』補題
\ref{sites-cohomology-lemma-dual-perfect-complex} を参照されたい。
$M \in D_\QCoh(\mathcal{O}_X)$ に対して
\begin{align*}
\Hom_{D(\mathcal{O}_Y)}(Rf_*M, K \otimes^\mathbf{L}_{\mathcal{O}_Y} L)
& =
\Hom_{D(\mathcal{O}_Y)}(
Rf_*M \otimes^\mathbf{L}_{\mathcal{O}_Y} K^\vee, L) \\
& =
\Hom_{D(\mathcal{O}_X)}(
M \otimes^\mathbf{L}_{\mathcal{O}_X} Lf^*K^\vee, a(L)) \\
& =
\Hom_{D(\mathcal{O}_X)}(M,
Lf^*K \otimes^\mathbf{L}_{\mathcal{O}_X} a(L))
\end{align*}
である。二番目の等号は $a$ の定義と射影公式
（『サイト上のコホモロジー』補題
\ref{sites-cohomology-lemma-projection-formula}）、またはより一般的な
『空間の導来圏』補題
\ref{spaces-perfect-lemma-cohomology-base-change} による。% P10-E0270
したがって Yoneda の補題により主張が従う。
\end{proof}

\begin{lemma}
\label{spaces-duality-lemma-restriction-compare-with-pullback}
図式 (\ref{spaces-duality-equation-base-change}) が与えられているとする。
$K \in D_\QCoh(\mathcal{O}_Y)$ とする。図式
$$
\xymatrix{
L(g')^*(Lf^*K \otimes^\mathbf{L}_{\mathcal{O}_X} a(\mathcal{O}_Y))
\ar[r] \ar[d] & L(g')^*a(K) \ar[d] \\
L(f')^*Lg^*K \otimes_{\mathcal{O}_{X'}}^\mathbf{L} a'(\mathcal{O}_{Y'})
\ar[r] & a'(Lg^*K)
}
$$
は可換である。ここで水平矢印は $K$ と $Lg^*K$ に対する写像
(\ref{spaces-duality-equation-compare-with-pullback}) であり、垂直写像は
『サイト上のコホモロジー』評注
\ref{sites-cohomology-remark-base-change} および
(\ref{spaces-duality-equation-base-change-map}) を用いて構成される。
\end{lemma}

\begin{proof}
この証明では $Rf_*$ の代わりに $f_*$、$Lf^*$ の代わりに
$f^*$ などと書き、$\otimes^\mathbf{L}_{\mathcal{O}_X}$ などの代わりに
$\otimes$ と書く。(\ref{spaces-duality-equation-compare-with-pullback}) を合成
\begin{align*}
f^*K \otimes a(\mathcal{O}_Y)
& \to
a(f_*(f^*K \otimes a(\mathcal{O}_Y))) \\
& \leftarrow
% P10-E0271
a(K \otimes f_*a(\mathcal{O}_K)) \\
& \to
a(K \otimes \mathcal{O}_Y) \\
& \to
a(K)
\end{align*}
として書く。ここで最初の矢印は単位 $\eta_f$、二番目の矢印は
『サイト上のコホモロジー』式
(\ref{sites-cohomology-equation-projection-formula-map}) に $a$ を適用したもの
であり、これは『空間の導来圏』補題
\ref{spaces-perfect-lemma-cohomology-base-change} により同型である。
三番目の矢印は $\text{id}_K \otimes \text{Tr}_f$ に $a$ を適用したもの、
四番目の矢印は同型 $K \otimes \mathcal{O}_Y = K$ に $a$ を適用したものである。
補題の証明は、これらの写像がそれぞれ補題の主張にある可換正方形を
生じさせることを示すことからなる。$\eta_f$ と $\text{Tr}_f$ については、
これは補題 \ref{spaces-duality-lemma-unit-and-base-change} および
\ref{spaces-duality-lemma-trace-map-and-base-change} である。
『サイト上のコホモロジー』式
(\ref{sites-cohomology-equation-projection-formula-map}) を用いる矢印については、
『サイト上のコホモロジー』評注
\ref{sites-cohomology-remark-compatible-with-diagram} である。
乗法写像については明らかである。これで証明は完了する。
\end{proof}




\section{固有平坦射に対する順像の右随伴}
\label{spaces-duality-section-proper-flat}

\noindent
準コンパクトかつ準分離な代数空間の間の、固有、平坦かつ有限表示な射に対して、
順像の右随伴はいくつかの顕著な性質をもつ。

\begin{lemma}
\label{spaces-duality-lemma-proper-flat}
$S$ をスキームとする。
$Y$ を $S$ 上の準コンパクトかつ準分離な代数空間とする。
$f : X \to Y$ を、固有、平坦かつ有限表示な代数空間の射とする。
$a$ を、補題 \ref{spaces-duality-lemma-twisted-inverse-image} の
$Rf_* : D_\QCoh(\mathcal{O}_X) \to D_\QCoh(\mathcal{O}_Y)$
の右随伴とする。このとき $a$ は直和と可換である。
\end{lemma}

\begin{proof}
$P$ を $D(\mathcal{O}_X)$ の完全対象とする。『空間の導来圏』補題
\ref{spaces-perfect-lemma-flat-proper-perfect-direct-image-general}
により、複体 $Rf_*P$ は $Y$ 上完全である。
$K_i$ を $D_\QCoh(\mathcal{O}_Y)$ の対象の族とする。このとき
\begin{align*}
\Hom_{D(\mathcal{O}_X)}(P, a(\bigoplus K_i))
& =
\Hom_{D(\mathcal{O}_Y)}(Rf_*P, \bigoplus K_i) \\
& =
\bigoplus \Hom_{D(\mathcal{O}_Y)}(Rf_*P, K_i) \\
& =
\bigoplus \Hom_{D(\mathcal{O}_X)}(P, a(K_i))
\end{align*}
である。これは完全対象がコンパクトであることによる
（『空間の導来圏』命題
\ref{spaces-perfect-proposition-compact-is-perfect}）。
$D_\QCoh(\mathcal{O}_X)$ は完全生成対象をもつ
（『空間の導来圏』定理
\ref{spaces-perfect-theorem-bondal-van-den-Bergh}）ので、写像
$\bigoplus a(K_i) \to a(\bigoplus K_i)$ は同型である。
すなわち $a$ は直和と可換である。
\end{proof}

\begin{lemma}
\label{spaces-duality-lemma-compare-with-pullback-flat-proper}
$S$ をスキームとする。
$Y$ を $S$ 上の準コンパクトかつ準分離な代数空間とする。
$f : X \to Y$ を、固有、平坦かつ有限表示な代数空間の射とする。
写像 (\ref{spaces-duality-equation-compare-with-pullback}) は
$D_\QCoh(\mathcal{O}_Y)$ のすべての対象 $K$ に対して同型である。
\end{lemma}

\begin{proof}
補題 \ref{spaces-duality-lemma-proper-flat} により、$a$ は直和と可換する。
したがって、(\ref{spaces-duality-equation-compare-with-pullback}) が同型となる
$D_\QCoh(\mathcal{O}_Y)$ の対象全体は、厳密充満、飽和、三角な
$D_\QCoh(\mathcal{O}_Y)$ の部分圏であり、さらに直和をとる操作で保たれる。
$D_\QCoh(\mathcal{O}_Y)$ は加群圏であり
（『空間の導来圏』定理
\ref{spaces-perfect-theorem-DQCoh-is-Ddga}）、一つの完全対象により生成される
（『空間の導来圏』定理
\ref{spaces-perfect-theorem-bondal-van-den-Bergh}）。
したがって『代数詳論』評注
\ref{more-algebra-remark-P-resolution} と同様に論じることで、
(\ref{spaces-duality-equation-compare-with-pullback}) が一つの完全対象に対して同型である
ことを示せば十分である。しかし、完全対象に対する結果は補題
\ref{spaces-duality-lemma-compare-with-pullback-perfect} で既に成り立つ。
\end{proof}

\begin{lemma}
\label{spaces-duality-lemma-properties-relative-dualizing}
$Y$ をアフィン・スキームとする。
$f : X \to Y$ を、固有、平坦かつ有限表示な代数空間の射とする。
$a$ を、補題 \ref{spaces-duality-lemma-twisted-inverse-image} の
$Rf_* : D_\QCoh(\mathcal{O}_X) \to D_\QCoh(\mathcal{O}_Y)$
の右随伴とする。このとき
\begin{enumerate}
\item $a(\mathcal{O}_Y)$ は $D(\mathcal{O}_X)$ の $Y$-完全対象である。
\item $Rf_*a(\mathcal{O}_Y)$ の正次数のコホモロジー層は消える。
\item $\mathcal{O}_X \to
R\SheafHom_{\mathcal{O}_X}(a(\mathcal{O}_Y), a(\mathcal{O}_Y))$
は同型である。
\end{enumerate}
\end{lemma}

\begin{proof}
% P10-E0272
$D(\mathcal{O}_X)$ の完全対象 $E$ に対して
\begin{align*}
Rf_*(E \otimes_{\mathcal{O}_X}^\mathbf{L} \omega_{X/Y}^\bullet)
& =
Rf_*R\SheafHom_{\mathcal{O}_X}(E^\vee, \omega_{X/Y}^\bullet) \\
& =
R\SheafHom_{\mathcal{O}_Y}(Rf_*E^\vee, \mathcal{O}_Y) \\
& =
(Rf_*E^\vee)^\vee
\end{align*}
である。最初の等号については『サイト上のコホモロジー』補題
\ref{sites-cohomology-lemma-dual-perfect-complex} を参照されたい。
二番目の等号については、補題 \ref{spaces-duality-lemma-iso-on-RSheafHom}、
評注 \ref{spaces-duality-remark-iso-on-RSheafHom}、および『空間の導来圏』補題
\ref{spaces-perfect-lemma-flat-proper-perfect-direct-image-general}
を参照されたい。三番目の等号は双対の定義である。
特に、これらの参照結果は、得られた対象が $D(\mathcal{O}_Y)$ の
完全対象であることも示す。したがって『空間の射詳論』補題
\ref{spaces-more-morphisms-lemma-characterize-relatively-perfect}
により、$\omega_{X/Y}^\bullet$ は $Y$-完全である。これで (1) が示された。

\medskip\noindent
$M$ を $D_\QCoh(\mathcal{O}_Y)$ の対象とする。このとき
\begin{align*}
\Hom_Y(M, Rf_*a(\mathcal{O}_Y)) & =
\Hom_X(Lf^*M, a(\mathcal{O}_Y)) \\
& =
\Hom_Y(Rf_*Lf^*M, \mathcal{O}_Y) \\
& =
% P10-E0274
\Hom_Y(M \otimes_{\mathcal{O}_Y}^\mathbf{L} Rf_*\mathcal{O}_Y, \mathcal{O}_Y)
\end{align*}
である。最初の等号は『サイト上のコホモロジー』補題
\ref{sites-cohomology-lemma-adjoint} による。
二番目の等号は $a$ の構成による。三番目の等号は『空間の導来圏』補題
\ref{spaces-perfect-lemma-cohomology-base-change} による。% P10-E0273
ある $N$ に対して $Rf_*\mathcal{O}_X$ は Tor 振幅 $[0, N]$ をもつ
完全対象であることを思い出そう。『空間の導来圏』補題
\ref{spaces-perfect-lemma-flat-proper-perfect-direct-image-general}
を参照されたい。したがって $Rf_*\mathcal{O}_X$ は、次数 $[0,N]$ に
置かれた有限射影加群の複体で表せる
（『代数詳論』補題 \ref{more-algebra-lemma-perfect} と
$Y$ がアフィンであることを用いる）。ゆえに、ある $i > 0$ に対して
$M = \mathcal{O}_Y[-i]$ なら、最後の群は零である。
$Y$ はアフィンなので、$i > 0$ に対して
$H^i(Rf_*a(\mathcal{O}_Y)) = 0$ である。これで (2) が示された。

\medskip\noindent
$E$ を $D_\QCoh(\mathcal{O}_X)$ の完全対象とする。このとき
\begin{align*}
% P10-E0275
\begin{aligned}
&\Hom_X(E, \\
&R\SheafHom_{\mathcal{O}_X}(a(\mathcal{O}_Y), a(\mathcal{O}_Y))
\end{aligned}
& =
\begin{aligned}
&\Hom_X(E \otimes_{\mathcal{O}_X}^\mathbf{L} a(\mathcal{O}_Y), \\
&a(\mathcal{O}_Y))
\end{aligned} \\
& =
\begin{aligned}
&\Hom_Y(Rf_*(E \otimes_{\mathcal{O}_X}^\mathbf{L} a(\mathcal{O}_Y)), \\
&\mathcal{O}_Y)
\end{aligned} \\
& =
\begin{aligned}
&\Hom_Y(Rf_*(R\SheafHom_{\mathcal{O}_X}(E^\vee, a(\mathcal{O}_Y))), \\
&\mathcal{O}_Y)
\end{aligned} \\
& =
\begin{aligned}
&\Hom_Y(R\SheafHom_{\mathcal{O}_Y}(Rf_*E^\vee, \mathcal{O}_Y), \\
&\mathcal{O}_Y)
\end{aligned} \\
& =
R\Gamma(Y, Rf_*E^\vee) \\
& =
\Hom_X(E, \mathcal{O}_X)
\end{align*}
である。最初の等号は『サイト上のコホモロジー』補題
\ref{sites-cohomology-lemma-internal-hom} による。
二番目の等号は $a$ の定義である。三番目の等号は完全双対複体
$E^\vee$ の構成から来る。『サイト上のコホモロジー』補題
\ref{sites-cohomology-lemma-dual-perfect-complex} を参照されたい。
四番目の等号は、証明の最初の段落で示した等式
$Rf_*R\SheafHom_{\mathcal{O}_X}(E^\vee, \omega_{X/Y}^\bullet) =
R\SheafHom_{\mathcal{O}_Y}(Rf_*E^\vee, \mathcal{O}_Y)$ から従う。
五番目の等号は完全複体の二重双対性
（『サイト上のコホモロジー』補題
\ref{sites-cohomology-lemma-dual-perfect-complex}）と、
『空間の導来圏』補題
\ref{spaces-perfect-lemma-flat-proper-perfect-direct-image-general}
によって $Rf_*E$ が完全であることによる。最後の等号は $f$ に対する
Leray である。この等式列は本質的に、Yoneda の補題により (3) が
成り立つことを示している。実際、『空間の導来圏』補題
\ref{spaces-perfect-lemma-quasi-coherence-internal-hom} により、対象
$R\SheafHom(a(\mathcal{O}_Y), a(\mathcal{O}_Y))$ は
$D_\QCoh(\mathcal{O}_X)$ に属する。上で $E = \mathcal{O}_X$ とすると、
$\text{id}_{\mathcal{O}_X} \in \Hom_X(\mathcal{O}_X, \mathcal{O}_X)$
に対応する写像
$\alpha : \mathcal{O}_X \to
R\SheafHom_{\mathcal{O}_X}(a(\mathcal{O}_Y), a(\mathcal{O}_Y))$
を得る。上のすべての同型は $E$ について関手的なので、
$\alpha$ の錐は $D_\QCoh(\mathcal{O}_X)$ の対象 $C$ であり、
すべての完全 $E$ に対して $\Hom(E, C) = 0$ を満たす。
完全対象は生成するので（『空間の導来圏』定理
\ref{spaces-perfect-theorem-bondal-van-den-Bergh}）、
$\alpha$ は同型である。
\end{proof}




\section{固有平坦射に対する相対双対化複体}
\label{spaces-duality-section-relative-dualizing-proper-flat}

\noindent
『スキームの双対性』第
\ref{duality-section-proper-flat} 節および
\ref{duality-section-relative-dualizing-complexes} 節と、
第 \ref{spaces-duality-section-proper-flat} 節の内容に動機づけられて、
次の定義を与える。

\begin{definition}
\label{spaces-duality-definition-relative-dualizing-proper-flat}
$S$ をスキームとする。$f : X \to Y$ を、$S$ 上の代数空間の間の
固有かつ平坦な有限表示射とする。
$X/Y$ の \emph{相対双対化複体}とは、
$D(\mathcal{O}_X)$ の $Y$-完全対象 $\omega_{X/Y}^\bullet$ と写像
$$
\tau : Rf_*\omega_{X/Y}^\bullet \longrightarrow \mathcal{O}_Y
$$
からなる対 $(\omega_{X/Y}^\bullet, \tau)$ であって、任意のデカルト正方形
$$
\xymatrix{
X' \ar[r]_{g'} \ar[d]_{f'} & X \ar[d]^f \\
Y' \ar[r]^g & Y
}
$$
に対し、$Y'$ がアフィンスキームならば、対
$(L(g')^*\omega_{X/Y}^\bullet, Lg^*\tau)$ が、第
\ref{spaces-duality-section-twisted-inverse-image}、
\ref{spaces-duality-section-base-change-map}、
\ref{spaces-duality-section-base-change-II}、
\ref{spaces-duality-section-trace}、
\ref{spaces-duality-section-compare-with-pullback}、および
\ref{spaces-duality-section-proper-flat} 節で考察した対
$(a'(\mathcal{O}_{Y'}), \text{Tr}_{f', \mathcal{O}_{Y'}})$
と同型になるものをいう。
\end{definition}

\noindent
ここでいくつか注意しておくべきことがある。
\begin{enumerate}
\item 定義 \ref{spaces-duality-definition-relative-dualizing-proper-flat} では、
$\omega_{X/Y}^\bullet$ が $Y$-完全であるという仮定を省いてよい。
実際、$Y'$ を $Y$ のアフィンによるエタール被覆の各要素にわたって動かすと、
補題 \ref{spaces-duality-lemma-properties-relative-dualizing} より、
$\omega_{X/Y}^\bullet$ の $X$ のあるエタール被覆の各要素への制限は
$Y$-完全である。したがって $\omega_{X/Y}^\bullet$ は $Y$-完全である。
『空間の射についてさらに』第
\ref{spaces-more-morphisms-section-relatively-perfect} 節を参照せよ。
\item 相対双対化複体 $(\omega_{X/Y}^\bullet, \tau)$ と、定義
\ref{spaces-duality-definition-relative-dualizing-proper-flat} にあるデカルト正方形を考える。
同型
$(L(g')^*\omega_{X/Y}^\bullet, Lg^*\tau) \cong
(a'(\mathcal{O}_{Y'}), \text{Tr}_{f', \mathcal{O}_{Y'}})$
の存在を、次のように捉えることにする。すなわち、任意の
$M' \in D_\QCoh(\mathcal{O}_{X'})$ に対して写像
$$
\Hom_{X'}(M', L(g')^*\omega_{X/Y}^\bullet)
\longrightarrow
\Hom_{Y'}(Rf'_*M', \mathcal{O}_{Y'}),\quad
\varphi' \longmapsto Lg^*\tau \circ Rf'_*\varphi'
$$
が同型であるということである。これは $a'$ の定義と、第
\ref{spaces-duality-section-trace} 節の議論から従う。特に、米田の補題により
この同型は一意である。
\item $Y$ 自身がアフィンならば、相対双対化複体
$(\omega_{X/Y}^\bullet, \tau)$ は存在し、
$(a(\mathcal{O}_Y), \text{Tr}_{f, \mathcal{O}_Y})$ と標準的に同型である。
ここで $a$ は補題 \ref{spaces-duality-lemma-twisted-inverse-image} にある $Rf_*$ の右随伴であり、
$\text{Tr}_f$ は第 \ref{spaces-duality-section-trace} 節にあるものである。
実際、定義にある図式を与えると、補題 \ref{spaces-duality-lemma-more-base-change} により
同型 $L(g')^*a(\mathcal{O}_Y) \to a'(\mathcal{O}_{Y'})$ を得るが、これは補題
\ref{spaces-duality-lemma-trace-map-and-base-change} によりトレース写像と両立する。
\end{enumerate}
これにより、局所的に与えられた相対双対化複体を大域的なものへ
貼り合わせるためにちょうど十分な情報が得られる。
以下の補題の証明は読み飛ばしてもよい。

\begin{lemma}
\label{spaces-duality-lemma-relative-dualizing-RHom}
% P10-E0278
$S$ をスキームとする。$X \to Y$ を有限表示な固有平坦射とする。
$(\omega_{X/Y}^\bullet, \tau)$ が相対双対化複体ならば、
$\mathcal{O}_X \to
R\SheafHom_{\mathcal{O}_X}(\omega_{X/Y}^\bullet, \omega_{X/Y}^\bullet)$
は同型であり、$Rf_*\omega_{X/Y}^\bullet$ の正次数のコホモロジー層は消える。
\end{lemma}

\begin{proof}
$Y$ 上エタールなアフィンスキームへの基底変換後にこれを証明すれば十分であり、
その場合は補題 \ref{spaces-duality-lemma-properties-relative-dualizing} から従う。
\end{proof}

\begin{lemma}
\label{spaces-duality-lemma-uniqueness-relative-dualizing}
$S$ をスキームとする。$X \to Y$ を有限表示な固有平坦射とする。
$(\omega_j^\bullet, \tau_j)$、$j = 1, 2$ が $X/Y$ 上の二つの
相対双対化複体ならば、一意な同型
$(\omega_1^\bullet, \tau_1) \to (\omega_2^\bullet, \tau_2)$ が存在する。
\end{lemma}

\begin{proof}
$g : Y' \to Y$ を、$Y'$ がアフィンスキームであるようなエタール射とし、
$X' = Y' \times_Y X$ を基底変換とする。
定義 \ref{spaces-duality-definition-relative-dualizing-proper-flat} とその後の議論により、
一意な同型
$\iota : (\omega_1^\bullet|_{X'}, \tau_1|_{Y'}) \to
(\omega_2^\bullet|_{X'}, \tau_2|_{Y'})$ が存在する。
$Y'' \to Y'$ がアフィンのさらなるエタール射で
$X'' = Y'' \times_Y X$ ならば、$\iota|_{X''}$ は一意な同型
$(\omega_1^\bullet|_{X''}, \tau_1|_{Y''}) \to
(\omega_2^\bullet|_{X''}, \tau_2|_{Y''})$ である（一意性による）。
また、補題 \ref{spaces-duality-lemma-relative-dualizing-RHom} により
$$
\text{Ext}^p_{X'}(\omega_1^\bullet|_{X'}, \omega_2^\bullet|_{X'}) = 0,
\quad p < 0
$$
である。実際、
$\mathcal{O}_{X'} \cong
R\SheafHom_{\mathcal{O}_{X'}}(\omega_1^\bullet|_{X'}, \omega_1^\bullet|_{X'})
\cong
R\SheafHom_{\mathcal{O}_{X'}}(\omega_1^\bullet|_{X'}, \omega_2^\bullet|_{X'})$
である。

\medskip\noindent
各 $V_n = \coprod_{i \in I_n} Y_{n, i}$ において $Y_{n, i}$ がアフィンとなるような
エタール超被覆 $b : V \to Y$ を選ぶ。
これは『超被覆』補題 \ref{hypercovering-lemma-hypercovering-object} と評注
\ref{hypercovering-remark-take-unions-hypercovering-X} により可能である
（補題で得られた超被覆を、各次数で非交和をもつものに置き換える）。
$X_{n, i} = Y_{n, i} \times_Y X$ および $U_n = V_n \times_Y X$ とおくと、
エタール超被覆 $a : U \to X$ を得る（『超被覆』補題
\ref{hypercovering-lemma-hypercovering-morphism-sites}）。ここで
$U_n = \coprod X_{n, i}$ である。
$a : U \to X$ と複体 $\omega_1^\bullet$、$\omega_2^\bullet$ に対して、
『単体的空間』補題
\ref{spaces-simplicial-lemma-fppf-neg-ext-zero-hom} の仮定が満たされる。
したがって一意な射 $\iota : \omega_1^\bullet \to \omega_2^\bullet$ を得て、
その $X_{0, i}$ への制限は一意な同型
$(\omega_1^\bullet|_{X_{0, i}}, \tau_1|_{Y_{0, i}}) \to
(\omega_2^\bullet|_{X_{0, i}}, \tau_2|_{Y_{0, i}})$
である。% P10-E0280
さらに、図式
$$
\xymatrix{
% P10-E0279
Rf_*\omega_1^\bullet \ar[rd]_{\tau_1} \ar[rr]_{Rf_*\iota} & &
Rf_*\omega_1^\bullet \ar[ld]^{\tau_2} \\
& \mathcal{O}_Y
}
$$
が可換であることを示さなければならない。しかし、補題
\ref{spaces-duality-lemma-relative-dualizing-RHom} により
$Rf_*\omega_1^\bullet$ と $Rf_*\omega_2^\bullet$ の正次数のコホモロジー層は消える。
したがって、この可換性はアフィン $Y_{0, i}$ への制限後に証明すればよく、
そこでは構成により成り立つ。
\end{proof}

\begin{lemma}
\label{spaces-duality-lemma-covering-enough}
% P10-E0281
$S$ をスキームとする。$X \to Y$ を有限表示な固有平坦射とする。
$(\omega^\bullet, \tau)$ を、$D(\mathcal{O}_X)$ の $Y$-完全対象と写像
$\tau : Rf_*\omega^\bullet \to \mathcal{O}_Y$ からなる対とする。
デカルト図式
$$
\xymatrix{
X_i \ar[r]_{g_i'} \ar[d]_{f_i} & X \ar[d]^f \\
Y_i \ar[r]^{g_i} & Y
}
$$
があり、$Y_i$ はアフィンで、$\{g_i : Y_i \to Y\}$ はエタール被覆であり、
さらに定義 \ref{spaces-duality-definition-relative-dualizing-proper-flat} にあるような対の同型
$(\omega^\bullet|_{X_i}, \tau|_{Y_i})
\to (a_i(\mathcal{O}_{Y_i}), \text{Tr}_{f_i, \mathcal{O}_{Y_i}})$
があると仮定する。このとき $(\omega^\bullet, \tau)$ は $Y$ 上の $X$ の
相対双対化複体である。
\end{lemma}

\begin{proof}
$g : Y' \to Y$ および $X', f', g', a'$ を、定義
\ref{spaces-duality-definition-relative-dualizing-proper-flat} にあるものとする。
$((\omega')^\bullet, \tau') = (L(g')^*\omega^\bullet, Lg^*\tau)$ とおく。
$\{Y_i \times_Y Y' \to Y'\}$ を細分する、アフィンによる有限エタール被覆
$\{Y'_j \to Y'\}$ を取ることができる（『位相』補題
\ref{topologies-lemma-etale-affine}）。
したがって各 $j$ に対し、ある $i_j$ と $Y$ 上の射
$k_j : Y'_j \to Y_{i_j}$ が存在する。ファイバー積
$$
\xymatrix{
X'_j \ar[r]_{h_j'} \ar[d]_{f'_j} &
X' \ar[d]^{f'} \\
Y'_j \ar[r]^{h_j} & Y'
}
$$
を考える。$k_j$ の $f_{i_j}$ による基底変換として誘導される射を
$k'_j : X'_j \to X_{i_j}$ と書く。与えられた同型を射 $k'_j$ によって
$Y'_j$ へ制限すると、対の同型
% P10-E0282
$((\omega')^\bullet|_{X'_j}, \tau'|_{Y'_j})
\to (a_j(\mathcal{O}_{Y'_j}), \text{Tr}_{f'_j, \mathcal{O}_{Y'_j}})$
を得る。$f : X \to Y$ を $f' : X' \to Y'$ で置き換えれば、
次の段落で解く問題に帰着する。

\medskip\noindent
$Y$ がアフィンであると仮定する。問題は $(\omega^\bullet, \tau)$ が
次の対と同型であることを示すことである。
\par\noindent
$\begin{aligned}
&(\omega_{X/Y}^\bullet, \text{Tr}) \\
&= (a(\mathcal{O}_Y), \text{Tr}_{f, \mathcal{O}_Y})
\end{aligned}$
。
被覆 $\{Y_i \to Y\}$ は、アフィンの単一の全射エタール射
$\{g : Y' \to Y\}$ で与えられていると仮定してよい。
実際、まず $\{g_i: Y_i \to Y\}$ を有限部分被覆で置き換え、次いで
$g = \coprod g_i :  Y' = \coprod Y_i \to Y$ とおけばよい。細部は省略する。
$X' = Y' \times_Y X$ とおき、$f', g'$ を定義
\ref{spaces-duality-definition-relative-dualizing-proper-flat} にある写像とする。
このとき、与えられているのは同型
$$
(\omega^\bullet|_{X'}, \tau|_{Y'}) \to
(a'(\mathcal{O}_{Y'}), \text{Tr}_{f', \mathcal{O}_{Y'}})
$$
だけである。$(\omega_{X/Y}^\bullet, \text{Tr})$ は相対双対化複体であるから
（定義 \ref{spaces-duality-definition-relative-dualizing-proper-flat} の後の議論を参照）、
一意な同型
$$
(\omega_{X/Y}^\bullet|_{X'}, \text{Tr}|_{Y'}) \to
(a'(\mathcal{O}_{Y'}), \text{Tr}_{f', \mathcal{O}_{Y'}})
$$
が存在する。例えば補題 \ref{spaces-duality-lemma-uniqueness-relative-dualizing} により一意である。
二つの表示された同型を組み合わせると、同型
$$
\alpha :
(\omega^\bullet|_{X'}, \tau|_{Y'}) \to
(\omega_{X/Y}^\bullet|_{X'}, \text{Tr}|_{Y'})
$$
を得る。$Y'' = Y' \times_Y Y'$ および $X'' = Y'' \times_Y X$ とおく。% P10-E0283
$\alpha$ の $X''$ への二つの引戻しは、一意性により再び一致しなければならない。
$X'$ 上で $\omega_{X'/Y'}^\bullet$ の負次数の自己 Ext が消え
（補題 \ref{spaces-duality-lemma-relative-dualizing-RHom}）、しかも任意の射影
$Y' \times_Y \ldots \times_Y Y' \to Y'$ による引戻し後にもこのことが成り立つので
（細部は省略する。補題 \ref{spaces-duality-lemma-uniqueness-relative-dualizing} の証明と比較せよ）、
『単体的空間』補題 \ref{spaces-simplicial-lemma-fppf-neg-ext-zero-hom} により、
$\alpha$ は $X$ 上の同型 $\omega^\bullet \to \omega_{X/Y}^\bullet$ に降下する。
\end{proof}

\begin{lemma}
\label{spaces-duality-lemma-existence-relative-dualizing}
% P10-E0284
$S$ をスキームとする。$X \to Y$ を有限表示な固有平坦射とする。
相対双対化複体 $(\omega_{X/Y}^\bullet, \tau)$ が存在する。
\end{lemma}

\begin{proof}
各 $V_n = \coprod_{i \in I_n} Y_{n, i}$ において $Y_{n, i}$ がアフィンとなるような
エタール超被覆 $b : V \to Y$ を選ぶ。
これは『超被覆』補題 \ref{hypercovering-lemma-hypercovering-object} と評注
\ref{hypercovering-remark-take-unions-hypercovering-X} により可能である
（補題で得られた超被覆を、各次数で非交和をもつものに置き換える）。
$X_{n, i} = Y_{n, i} \times_Y X$ および $U_n = V_n \times_Y X$ とおくと、
エタール超被覆 $a : U \to X$ を得る（『超被覆』補題
\ref{hypercovering-lemma-hypercovering-morphism-sites}）。ここで
$U_n = \coprod X_{n, i}$ である。
各 $n, i$ に対し、$X_{n, i}/Y_{n, i}$ 上の相対双対化複体
$(\omega_{n, i}^\bullet, \tau_{n, i})$ が存在する。定義
\ref{spaces-duality-definition-relative-dualizing-proper-flat} の後の議論を参照せよ。
$\varphi : [m] \to [n]$ と $i \in I_n$ に対し、
% P10-E0285
$g_{\varphi, i} : Y_{n, i} \to Y_{m, \alpha(\varphi)}$ および
$g'_{\varphi, i} : X_{n, i} \to X_{m, \alpha(\varphi)}$ を、
与えられた超被覆の構造の一部をなす射とする
（『超被覆』第 \ref{hypercovering-section-hypercovering-sites} 節）。
このとき一意な同型 % P10-E0286
$$
% P10-E0287
\iota_{n, i, \varphi} :
(L(g'_{n, i})^*\omega_{n, i}^\bullet, Lg_{n, i}^*\tau_{n, i})
\longrightarrow
(\omega_{m, \alpha(\varphi)(i)}^\bullet, \tau_{m, \alpha(\varphi)(i)})
$$
が存在する。定義 \ref{spaces-duality-definition-relative-dualizing-proper-flat} の後の議論を参照せよ。
補題 \ref{spaces-duality-lemma-relative-dualizing-RHom} により、
$X_{n, i}$ 上で $\omega_{n, i}^\bullet$ の負次数の自己 Ext は消える。
$U_n/V_n$ 上で、$i \in I_n$ に対する対
$(\omega_{n, i}^\bullet, \tau_{n, i})$ から構成される対を
$(\omega_n^\bullet, \tau_n)$ と書く。
$\varphi : [m] \to [n]$ と $i \in I_n$ に対し、
$g_\varphi : V_n \to V_m$ および $g'_\varphi : U_n \to U_m$ を、
単体的代数空間 $V$ と $U$ の構造の一部をなす射とする。
このとき、各成分上の同型から構成される対の一意な同型
$$
\iota_\varphi :
(L(g'_\varphi)^*\omega_n^\bullet, Lg_\varphi^*\tau_n)
\longrightarrow
(\omega_m^\bullet, \tau_m)
$$
が存在する。一意性により、これらの同型は『単体的空間』定義
\ref{spaces-simplicial-definition-cartesian-derived-modules}
で定式化された推移律条件を満たす。
『単体的空間』補題
\ref{spaces-simplicial-lemma-fppf-glue-neg-ext-zero} の仮定は、
$a : U \to X$、複体 $\omega_n^\bullet$、および同型 $\iota_\varphi$ に対して
満たされる\footnote{この補題で用いるのは $\omega_0^\bullet$ と二つの写像
% P10-E0288
$\delta_1^1, \delta_0^1 : [1] \to [0]$ だけである。
ここでは実際、加群の導来圏の単体的系がすでに与えられているので、
参照した補題の証明の最初の数行は読み飛ばしてよい。}。
したがって $D_\QCoh(\mathcal{O}_X)$ の対象 $\omega^\bullet$ と、
二つの同型 $\iota_{\delta^1_0}$、$\iota_{\delta^1_1}$ と両立する同型
$\iota_0 : \omega^\bullet|_{U_0} \to \omega_0^\bullet$ を得る。
最後に、『単体的空間』補題
\ref{spaces-simplicial-lemma-fppf-neg-ext-zero-hom} を適用して、
一意な射
$$
\tau : Rf_*\omega^\bullet \longrightarrow \mathcal{O}_Y
$$
を得る。その $V_0$ への制限は $\tau_0$ と一致する。
細部は省略する。例えば必要な負次数 Ext の消滅が得られる理由については、
補題 \ref{spaces-duality-lemma-uniqueness-relative-dualizing} の証明の末尾と比較せよ。% P10-E0289
補題 \ref{spaces-duality-lemma-covering-enough} により、対 $(\omega^\bullet, \tau)$ は
相対双対化複体であり、証明が完了する。
\end{proof}

\begin{lemma}
\label{spaces-duality-lemma-base-change-relative-dualizing}
$S$ をスキームとする。$S$ 上の代数空間のデカルト正方形
$$
\xymatrix{
X' \ar[d]_{f'} \ar[r]_{g'} & X \ar[d]^f \\
Y' \ar[r]^g & Y
}
$$
を考える。$X \to Y$ は固有、平坦かつ有限表示であると仮定する。
$(\omega_{X/Y}^\bullet, \tau)$ を $f$ の相対双対化複体とする。
このとき $(L(g')^*\omega_{X/Y}^\bullet, Lg^*\tau)$ は $f'$ の
相対双対化複体である。
\end{lemma}

\begin{proof}
『空間の射についてさらに』補題
\ref{spaces-more-morphisms-lemma-base-change-relatively-perfect} により、
$L(g')^*\omega_{X/Y}^\bullet$ は $Y'$-完全である。
定義 \ref{spaces-duality-definition-relative-dualizing-proper-flat} のもう一つの条件は、
ファイバー積の推移律から成り立つ。
\end{proof}







\section{スキームの場合との比較}
\label{spaces-duality-section-comparison}

\noindent
この節には、なお多くの内容を加えるべきである。

\begin{lemma}
\label{spaces-duality-lemma-compare}
$S$ をスキームとする。$f : X \to Y$ を、$S$ 上の準コンパクトかつ
準分離な代数空間の射とする。
$X$ と $Y$ は表現可能であると仮定し、$f$ を表現するスキームの射を
$f_0 : X_0 \to Y_0$ とする（不格好だが一時的な記法である）。
$a : D_\QCoh(\mathcal{O}_Y) \to D_\QCoh(\mathcal{O}_X)$ を、補題
\ref{spaces-duality-lemma-twisted-inverse-image} にある $Rf_*$ の右随伴とする。
$a_0 : D_\QCoh(\mathcal{O}_{Y_0}) \to D_\QCoh(\mathcal{O}_{X_0})$
を、『スキームの双対性』補題
\ref{duality-lemma-twisted-inverse-image} にある $Rf_*$ の右随伴とする。
このとき
$$
\xymatrix{
D_\QCoh(\mathcal{O}_{X_0})
\ar@{=}[rrrrrr]_{\text{『空間の導来圏』補題
\ref{spaces-perfect-lemma-derived-quasi-coherent-small-etale-site}}}
& & & & & &
D_\QCoh(\mathcal{O}_X) \\
D_\QCoh(\mathcal{O}_{Y_0}) \ar[u]^{a_0}
\ar@{=}[rrrrrr]^{\text{『空間の導来圏』補題
\ref{spaces-perfect-lemma-derived-quasi-coherent-small-etale-site}}}
& & & & & &
D_\QCoh(\mathcal{O}_Y) \ar[u]_a
}
$$
は可換である。
\end{lemma}

\begin{proof}
随伴の一意性と、『空間の導来圏』評注
\ref{spaces-perfect-remark-match-total-direct-images} の両立性から従う。
\end{proof}
